\documentclass[journal, a4paper]{IEEEtran}
%\documentclass[12pt,onecolumn]{IEEEtran}

%\usepackage{cite}
\usepackage{booktabs}
\usepackage{graphicx}
%\usepackage{psfrag}
%\usepackage{subfigure}
%\usepackage{hyperref}
\usepackage{stfloats}
\usepackage{placeins}
\usepackage{url}
\usepackage{cuted}
\usepackage{lipsum}
\usepackage{amsmath}
\usepackage{amssymb}
\usepackage{siunitx}
\sisetup{output-exponent-marker=\ensuremath{\mathrm{e}}}

\begin{document}

\title{Time-Optimal Heliocentric Transfers With a Constant-Power, Variable-$I_{\!sp}$ Engine}

\author{Jan~Ol\v{s}ina\\[1ex] Independent Researcher, Prague, Czech Republic\\
email: \texttt{jan.olsina@gmail.com},  ORCID: \texttt{0009-0003-9816-3958}.}


\maketitle

%\begin{strip}
\begin{abstract}
We study minimum-time heliocentric transfers for a spacecraft propelled by an electric thruster that draws constant electrical power $P$ while continuously varying its exhaust speed (variable $I_{\!sp}$). The vehicle is assumed to depart from and arrive on circular heliocentric orbits (i.e., initial and final velocities match the local circular velocity at the respective radii). First, we derive an analytic solution of the one-dimensional, gravity-free brachistochrone and discuss how a finite exhaust-speed ceiling modifies the solution, producing a boost–coast–brake structure. Next, we formulate the full planar Sun-field optimal-control problem, derive two closed-form first integrals, and show that the indirect formulation reduces to a seven-dimensional boundary-value problem. Finally, we present a practical numerical continuation strategy that obtains a coarse feasible endpoint via global optimization and then refines it by homotopy and Powell’s local solver. Numerical examples for a \SI{1}{GW} engine with an initial/dry mass of \SI{3000}{t}\,\(\to\)\,\SI{1000}{t} demonstrate Earth–Jupiter-class transfers in roughly $200$–$220$ days that commonly exploit a solar Oberth pass. 
\end{abstract}
%\end{strip}


\section{Introduction}
While electric propulsion (EP) is traditionally valued for its propellant efficiency, future ambitious endeavors such as interplanetary logistics networks or rapid sample return missions will demand not just efficiency but \emph{speed} \cite{Conway2010, PinoHowell2020, Fogel2020}. For such missions, the design focus shifts from optimizing propellant mass to minimizing flight time, making the spacecraft's power system the primary limiting factor \cite{Mazouffre2016}.
Reducing transfer times from years to months propels interest in high-power propulsion systems, with gigawatt-class architectures representing a plausible long-term goal. Several pathways to this power scale are being explored, including advanced Nuclear Electric Propulsion (NEP) \cite{Peakman2023, scott2019value}, next-generation deployable solar arrays \cite{ChangDiaz2019}, and beamed-energy concepts \cite{Sheerin2021}. A common feature of these technologies is the ability to provide nearly constant power over long thrust arcs.
This constant-power paradigm fundamentally alters the classic rocketry framework. Instead of the fixed exhaust velocity ($v_e$) of the Tsiolkovsky equation \cite{Tsiolkovsky1903}, the exhaust speed becomes a key control variable, allowing the engine to actively trade thrust for specific impulse. The Variable Specific Impulse Magnetoplasma Rocket (VASIMR) is an exemplary technology demonstrating this "constant-power throttling" over a wide operating range \cite{Bering2011}. Such throttleability is compelling for the fast transfers that motivate our study, as recent optimization work confirms that varying the specific impulse during flight can significantly improve trip time \cite{DiPasquale2024}.
In this paper, we develop the machinery to solve the time-optimal transfer problem in this high-power regime. We first derive an analytical solution for a 1D, gravity-free "brachistochrone" to build intuition. We then formulate the full planar Sun-field problem, derive two closed-form first integrals to simplify the dynamics, and present a robust numerical strategy to solve the resulting boundary-value problem. Our examples for a gigawatt-scale engine demonstrate Earth–Jupiter class transfers in a few months, frequently exploiting a solar Oberth maneuver.


\section{Variable-$I_{\!sp}$ Rocket Model}
Let $m(t)$ be the instantaneous mass ($m \ge m_\mathrm{dry}$) and $v_e(t)$ the exhaust speed. The propellant mass flow rate is the positive quantity $\dot m_f = -\dot m$. Neglecting plume inefficiencies, the thruster power $P$ (constant) and thrust magnitude $T$ are given by \cite{Sutton2017}:
\begin{equation}
  P = \tfrac12 (-\dot m)\,v_e^{2},\qquad
  T   = m\,a = (-\dot m)\,v_e,
  \label{eq:power_thrust}
\end{equation}
where $a$ is the resulting thrust-induced acceleration magnitude. Eliminating $v_e$ gives a direct relationship between the mass flow rate and the vehicle's acceleration:
\begin{equation}
  \dot m = -\frac{m^{2}a^{2}}{2P}.
  \label{eq:mass_flow}
\end{equation}
This equation forms the basis of our optimal control problem.

\section{Exact Solution for the Free-Space Case $(\mu=0)$}
\label{sec:free_space_solution}
With solar gravity absent ($\mu=0$), the translational degrees of freedom decouple. The problem simplifies to finding the minimum-time trajectory for a 1D drive from an initial state $(x, v) = (0, 0)$ at $t=0$ to a final state $(x, v) = (L, 0)$ at the final time $t_f$.

\subsection{The Variational Method and the Hamiltonian}
\label{sec:pmp_theory}
To find the optimal control history---in this case, the acceleration profile $a(t)$---we use the calculus of variations, specifically Pontryagin's Maximum Principle (PMP) \cite{Pontryagin1962}. This powerful method converts the problem of minimizing a cost functional (like total time) into a problem of maximizing a new function, the \emph{Hamiltonian}, at every point in time.

The Hamiltonian, $\mathcal{H}$, is constructed to package all essential information about the problem: the cost, the system dynamics, and the controls. The general formula is:
\begin{equation}
    \mathcal{H} = L + \boldsymbol{\lambda}^\top \mathbf{f}(\mathbf{x}, \mathbf{u}),
    \label{eq:hamiltonian_general}
\end{equation}
where:
\begin{itemize}
    \item $L$ is the \textbf{Lagrangian}, or the instantaneous cost. For a minimum-time problem, the total cost is $J = \int_0^{t_f} 1 \, dt$, so the instantaneous cost is simply $L=1$.
    \item $\mathbf{x}$ is the \textbf{state vector}, containing the variables that describe the system (e.g., position, velocity, mass).
    \item $\mathbf{u}$ is the \textbf{control vector}, containing the variables we can change to steer the system (e.g., acceleration).
    \item $\mathbf{f}(\mathbf{x}, \mathbf{u}) = \dot{\mathbf{x}}$ is the vector of \textbf{state dynamics}, describing how the system evolves.
    \item $\boldsymbol{\lambda}$ is the \textbf{costate vector}, a set of auxiliary variables introduced by the PMP. Intuitively, each costate $\lambda_i(t)$ represents the sensitivity of the final cost to a small change in the state $x_i$ at time $t$. It is the "price" of that state variable.
\end{itemize}

\subsection{Hamiltonian Construction for the Free-Space Case}
We apply the principle from Section~\ref{sec:pmp_theory}. First, we identify the components of the 1D problem:
\begin{itemize}
    \item \textbf{State vector:} ${\bf x} = (x, v, m)^\top$.
    \item \textbf{Control variable:} $u = a(t)$.
    \item \textbf{State dynamics} $\dot{\mathbf{x}} = \mathbf{f}(\mathbf{x}, u)$:
    \begin{align}
      \dot x &= v, \label{eq:state_x_free}\\
      \dot v &= a, \label{eq:state_v_free}\\
      \dot m &= -\frac{m^{2}a^{2}}{2P}. \label{eq:state_m_free}
    \end{align}
    \item \textbf{Costate vector:} $\boldsymbol{\lambda} = (\lambda_x, \lambda_v, \lambda_m)^\top$.
\end{itemize}
The term $\boldsymbol{\lambda}^\top \dot{\mathbf{x}}$ in the Hamiltonian expands to:
\[
    \boldsymbol{\lambda}^\top \dot{\mathbf{x}} = \lambda_x \dot{x} + \lambda_v \dot{v} + \lambda_m \dot{m} = \lambda_x v + \lambda_v a - \lambda_m \frac{m^2 a^2}{2P}.
\]
Adding the Lagrangian $L=1$, we obtain the full Hamiltonian $\mathcal{H}$:
\begin{equation}
    \mathcal{H} = 1 + \lambda_x v + \lambda_v a - \lambda_m \frac{m^2 a^2}{2P}.
\end{equation}
From this Hamiltonian, we derive the costate (adjoint) equations, given by $\dot{\boldsymbol{\lambda}} = -\partial\mathcal{H}/\partial\mathbf{x}$:
\begin{subequations} \label{eq:costate_free}
\begin{align}
  \dot\lambda_x &= -\frac{\partial\mathcal{H}}{\partial x} = 0, \label{eq:costate_x_free}\\
  \dot\lambda_v &= -\frac{\partial\mathcal{H}}{\partial v} = -\lambda_x, \label{eq:costate_v_free}\\
  \dot\lambda_m &= -\frac{\partial\mathcal{H}}{\partial m} = \lambda_m \frac{m a^2}{P}. \label{eq:costate_m_free}
\end{align}
\end{subequations}

\subsection{Optimal Control and First Integrals}
According to the PMP, the optimal control $a(t)$ must maximize the Hamiltonian at all times. We find it by setting $\partial\mathcal{H}/\partial a = 0$:
\begin{equation}
    \frac{\partial\mathcal{H}}{\partial a} = \lambda_v - \lambda_m \frac{m^2 a}{P} = 0 \implies a(t) = \frac{P \lambda_v(t)}{\lambda_m(t) m^2(t)}.
    \label{eq:a_of_lambda}
\end{equation}
From the costate dynamics, we can extract two integrals of motion. First, \eqref{eq:costate_x_free} directly implies:
\begin{equation}
    \lambda_x(t) = C_1 = \text{const}.
\end{equation}
Substituting this into \eqref{eq:costate_v_free} and integrating gives:
\begin{equation}
    \lambda_v(t) = -C_1 t + C_2,
\end{equation}
showing that the velocity costate is a linear function of time.

A second, less obvious, integral is found by examining the time derivative of the product $\lambda_m m^2$:
\begin{align}
    \frac{d}{dt}(\lambda_m m^2) &= \dot\lambda_m m^2 + \lambda_m (2m \dot m) \nonumber\\
    &= \left(\lambda_m \frac{m a^2}{P}\right) m^2 + \lambda_m \left(2m \left(-\frac{m^2 a^2}{2P}\right)\right) \nonumber\\
    &= \frac{\lambda_m m^3 a^2}{P} - \frac{\lambda_m m^3 a^2}{P} = 0.
\end{align}
Thus, the product $\lambda_m(t) m^2(t)$ is constant throughout the trajectory:
\begin{equation}
    \lambda_m(t) m^2(t) = C_m = \text{const}.
    \label{eq:mass_costate_integral}
\end{equation}
Substituting these integrals into the control law \eqref{eq:a_of_lambda}, we discover that the optimal acceleration profile must be linear in time:
\begin{equation}
    a(t) = \frac{P}{C_m}(-C_1 t + C_2) = K_1 t + K_0,
\end{equation}
where $K_1 = -PC_1/C_m$ and $K_0 = PC_2/C_m$ are constants.

\subsection{Deriving the Trajectory}
We now determine the constants $K_1$ and $K_0$ by applying the boundary conditions. Let the final time be $t_f$. First, we integrate $a(t)$ to find the velocity $v(t)$:
\begin{equation}
    v(t) = \int_0^t a(\tau) d\tau + v(0) = \tfrac{1}{2} K_1 t^2 + K_0 t.
\end{equation}
Applying the final velocity condition $v(t_f)=0$:
\begin{equation}
    \tfrac{1}{2} K_1 t_f^2 + K_0 t_f = 0 \implies K_0 = -\tfrac{1}{2} K_1 t_f.
\end{equation}
Substituting this back into $a(t)$ reveals its symmetric nature, a characteristic feature of such time-optimal "there-and-back" problems:
\begin{equation}
    a(t) = K_1 t - \tfrac{1}{2} K_1 t_f = K_1 (t - t_f/2).
\end{equation}
Next, we integrate $v(t)$ to find the position $x(t)$:
\begin{equation}
    x(t) = \int_0^t v(\tau) d\tau + x(0) = \tfrac{1}{6} K_1 t^3 + \tfrac{1}{2} K_0 t^2.
\end{equation}
Applying the final position condition $x(t_f)=L$:
\begin{equation}
    L = \tfrac{1}{6} K_1 t_f^3 + \tfrac{1}{2} K_0 t_f^2.
\end{equation}
Substituting $K_0 = -\frac{1}{2} K_1 t_f$ into the equation for $L$ allows us to solve for $K_1$:
\begin{align}
    L &= \tfrac{1}{6} K_1 t_f^3 + \tfrac{1}{2}(-\tfrac{1}{2}K_1 t_f)t_f^2 \nonumber\\
      &= (\tfrac{1}{6} - \tfrac{1}{4})K_1 t_f^3 = -\tfrac{1}{12} K_1 t_f^3.
\end{align}
This gives $K_1 = -12L/t_f^3$. Consequently, $K_0 = -\frac{1}{2}(-\frac{12L}{t_f^3})t_f = 6L/t_f^2$.
The optimal acceleration is therefore:
\begin{equation}
  a(t)=\frac{12L}{t_f^{3}}\Bigl(\frac{t_f}{2}-t\Bigr).
  \label{eq:a_linear}
\end{equation}
This corresponds to an initial acceleration phase, followed by a symmetric braking phase, with 
$a(t)$ crossing zero at mid-time. $t=t_f/2$. By substituting the constants back and simplifying, we obtain the full trajectory:
\begin{align}
  v(t) &= \frac{6Lt}{t_f^2}\Bigl(1-\frac{t}{t_f}\Bigr),
         \label{eq:v_free_corrected}\\
  x(t) &= L\Bigl[3\bigl(\tfrac{t}{t_f}\bigr)^{2}
                   -2\bigl(\tfrac{t}{t_f}\bigr)^{3}\Bigr],
         \label{eq:x_free}
\end{align}
which describes a cubic-in-time ``S-curve'' path.

\subsection{Mass Evolution and Minimum Time}
The mass evolution is governed by the separable differential equation \eqref{eq:state_m_free}. Integrating gives:
\begin{align}
    \int_{m_i}^{m(t)} \frac{dm}{m^2} &= -\int_0^t \frac{a(\tau)^2}{2P} d\tau \nonumber\\ 
    \implies \frac{1}{m_i} - \frac{1}{m(t)} &= -\frac{1}{2P} \int_0^t a(\tau)^2 d\tau.
\end{align}
The integral of $a(t)^2$ is:
\begin{align}
    \int_0^t a(\tau)^2 d\tau &= \int_0^t \left(\frac{12L}{t_f^3}\right)^2 \left(\frac{t_f}{2}-\tau\right)^2 d\tau \nonumber\\
    &= \frac{144L^2}{t_f^6} \left[-\frac{1}{3}\left(\frac{t_f}{2}-\tau\right)^3\right]_0^t \nonumber\\
    &= \frac{48L^2}{t_f^6} \left(\frac{t_f^3}{8} - \left(\frac{t_f}{2}-t\right)^3 \right).
\end{align}
Defining 
\begin{equation}
f(s)=4s^{3}- 6s^{2}+3s,\label{eq:fdef}
\end{equation}
the mass relation becomes:
\begin{equation}
  \frac{1}{m(t)}-\frac{1}{m_i}
      =\frac{6L^{2}f(t/t_f)}{P\,t_f^{3}}.
  \label{eq:m_free}
\end{equation}
The total propellant consumed is found by setting $t=t_f$ (or $s=1$). If the entire propellant reserve is used, we have $m(t_f)=m_\mathrm{dry}$, which defines the minimum possible flight-time for the given distance $L$ and mass fraction:
\begin{equation}
  \Delta\equiv\frac{1}{m_\mathrm{dry}}-\frac{1}{m_i}
\end{equation} 
calculated as
\begin{equation}
  \Delta = \frac{6L^{2}f(1)}{P\,T_{\min}^{3}} = \frac{6L^2}{P\,T_{\min}^3}.
\end{equation}
Solving for $T_{\min}$ yields the final result:
\begin{equation}
  T_{\min}= \left[\,
       \frac{6L^{2}}
             {P\Delta}
       \right]^{1/3}.
  \label{eq:Tmin}
\end{equation}

Equation~\eqref{eq:Tmin} treats the dry mass as prescribed. If the propulsion hardware mass is instead allowed to scale with power, the same free-space benchmark can be extended by a simple dry-mass closure; the resulting scaling law is derived in Appendix~\ref{sec:appendix_mass_closure}.

\section{Transfers with a Finite Exhaust-Speed Ceiling}
\label{sec:capped_ve}
The linear-in-time acceleration profile from Sec.~\ref{sec:free_space_solution} is feasible only if the engine can raise its exhaust speed indefinitely, because $v_e(t)=2P/\!\bigl(m(t)a(t)\bigr)$ diverges when $a(t)\!\to\!0$ at mid-time. A realistic thruster, however, is limited to a maximum effective exhaust speed $v_{e,\max}$. This path constraint forces the optimal trajectory to switch the engine \emph{off} once the required $v_e$ reaches the ceiling, producing the well-known \emph{boost-coast-brake} structure \cite{Fogel2020, pontani2021optimal}. 

This section derives the explicit solution for this case by adapting the unlimited solution. We introduce a coast phase and scale the acceleration profile to satisfy all physical constraints. The problem reduces to solving a single nonlinear equation for the optimal duration of the coast phase.

\subsection{Setup and Scaled Trajectory}
The trajectory consists of a burn, a coast, and a symmetric braking burn. We define a dimensionless \emph{switch fraction} $\xi\equiv t_c/t_f\in(0,\tfrac12)$, where $t_c$ is the time when the engine is first shut down and $t_f$ is the total transfer time.

We adapt the acceleration profile from the unlimited case, $a_\infty(t)$ from \eqref{eq:a_linear}, by scaling it with a constant throttle factor $k$. The powered phases now run from $t=0$ to $t_c$ and symmetrically at the end. The three unknowns ($k$, $t_f$, and $\xi$) are determined by the three physical constraints: total distance, total propellant, and the exhaust speed ceiling. The derivation uses the recurring polynomial $f(\xi)$  defined in \eqref{eq:fdef}, which is positive for $0<\xi<\tfrac12$.

\subsection{Derivation via Constraint Matching}
Since the Pontryagin conditions are preserved on the powered arcs, the scaled linear-in-time profile with an inserted coast is an admissible PMP candidate; enforcing the three global constraints selects the optimal member of this family.

\paragraph{Distance Closure}
The total distance traveled, $L$, is the sum of distances covered during the first burn, the coast, and the final brake. Integrating the velocity profile for the scaled acceleration over the three segments and setting the sum to $L$ uniquely determines the throttle scale $k$ as a function of $\xi$:
\begin{equation}
k = \frac{1}{2f(\xi)}.
\end{equation}

\paragraph{Propellant Closure}
The total change in the inverse of mass, $\Delta$, must be consumed over the two symmetric powered phases. Integrating the mass-flow equation \eqref{eq:mass_flow} using the scaled acceleration $a(t) = k\, a_\infty(t)$ gives:
\begin{equation}
\Delta = 2 \int_0^{t_c} \frac{a(t)^2}{2P} dt = \frac{k^2}{P} \int_0^{t_c} a_\infty(t)^2 dt.
\end{equation}
Solving this for the total time $t_f$ and substituting $k=1/(2f(\xi))$ yields the flight time as a function of the switch fraction:
\begin{equation}
t_f(\xi)=
\Biggl[
\frac{3\,L^{2}}{P\,\Delta\,f(\xi)}
\Biggr]^{\!\!1/3}.
\label{eq:T_of_xi_new}
\end{equation}

\paragraph{Exhaust Speed Ceiling}
The final constraint is that the exhaust speed must not exceed its ceiling, $v_e(t) \le v_{e,\max}$. The critical point is at the switch-off time $t=t_c$, where acceleration is lowest and mass has decreased. We impose the limit exactly at this point: $v_e(t_c) = 2P/(m_c a(t_c)) = v_{e,\max}$. The mass at this point, $m_c$, is the harmonic mean of the initial and dry masses, $m_c=2m_i m_\mathrm{dry}/(m_i+m_\mathrm{dry})$, since half the propellant is consumed. Substituting the expressions for $a(t_c)$, $k$, and $t_f(\xi)$ into this ceiling condition gives a single nonlinear equation for the optimal switch fraction~$\xi_\star$:
\begin{align}
F(\xi)&=0, \quad \text{where} \\
F(\xi)&\equiv
\bigl(\tfrac12-\xi\bigr)\,v_{e,\max}
-\frac{P^{1/3}L^{1/3}f(\xi)^{1/3}}
     {3^{1/3}\Delta^{2/3}\,m_c}.
\label{eq:root_equation_new}
\end{align}
The function $F(\xi)$ is strictly decreasing on $0<\xi<\tfrac12$, so if a root exists, it is unique.

\subsection{Final Solution and Limiting Cases}
The optimal transfer is found by solving the scalar equation $F(\xi)=0$ for the root $\xi_\star \in (0, 1/2)$ using any robust 1-D method. This value is then used in \eqref{eq:T_of_xi_new} to find the minimum transfer time.

\paragraph*{Detecting whether the ceiling is active}
If $F(\tfrac12)\ge0$, the ceiling is never touched. The optimal strategy is to never coast, and the minimum time is given by the unlimited-engine benchmark~\eqref{eq:Tmin}. Otherwise, a unique root $\xi_\star$ exists, and the ceiling constraint is active.

\paragraph*{Limiting Cases}
\begin{itemize}
\item \textbf{Unlimited engine.} As $v_{e,\max}\!\to\!\infty$, the root of $F(\xi)=0$ migrates to $\xi\!\to\!\tfrac12$. In this limit, $f(\xi)\to 1/2$ and the flight time~\eqref{eq:T_of_xi_new} collapses to the cubic root benchmark~\eqref{eq:Tmin}, recovering the solution from Sec.~\ref{sec:free_space_solution}.
\item \textbf{Vanishing ceiling.} When $v_{e,\max}$ approaches the minimum useful value, $F(\xi)$ forces $\xi\!\to\!0^{+}$. The two powered legs degenerate into instantaneous \emph{impulses} separated by a long ballistic cruise, as intuition demands.
\end{itemize}

\section{Full Planar Sun-Field Model}
\label{sec:SunModel}
We now reinstate solar gravity ($\mu=GM_\odot$) and adopt polar coordinates ${\bf x}=(r,\theta)$ in the ecliptic. The control is the \emph{planar thrust acceleration} ${\bf a}_{T}=(a_r,a_\theta)$ produced by the variable-$I_{\!sp}$ engine. The state vector and its dynamics are:
\begin{subequations}\label{eq:StateDyn}
\begin{align}
  \dot r       &= v_r, \\
  \dot\theta   &= \frac{v_\theta}{r}, \\
  \dot v_r     &= \frac{v_\theta^{2}}{r}-\frac{\mu}{r^{2}} + a_r, \\
  \dot v_\theta&= -\frac{v_r v_\theta}{r} + a_\theta, \\
  \dot m       &= -\frac{m^{2}}{2P}\bigl(a_r^{2}+a_\theta^{2}\bigr).
\end{align}
\end{subequations}

\subsection{Pontryagin Maximum Principle (PMP)}
We again apply the variational framework from Section~\ref{sec:pmp_theory}. For this time-optimal problem, the performance index is $J=\displaystyle\int_{0}^{t_f}\!{\rm d}t$, meaning the Lagrangian is $L=1$. The components for the Hamiltonian construction are now:
\begin{itemize}
    \item \textbf{State vector:} $\mathbf{x} = (r, \theta, v_r, v_\theta, m)^\top$.
    \item \textbf{Control vector:} $\mathbf{u} = \mathbf{a}_T = (a_r, a_\theta)^\top$.
    \item \textbf{Costate vector:} $\boldsymbol{\lambda} = (\lambda_r, \lambda_\theta, \lambda_{v_r}, \lambda_{v_\theta}, \lambda_m)^\top$.
\end{itemize}
\medskip
The Hamiltonian, $\mathcal{H} = 1 + \boldsymbol{\lambda}^\top \dot{\mathbf{x}}$, must be maximized by the optimal control at all times. The dot product term $\boldsymbol{\lambda}^\top \dot{\mathbf{x}}$ is constructed by multiplying each costate by its corresponding state derivative from \eqref{eq:StateDyn}:
\begin{align}
\boldsymbol{\lambda}^\top \dot{\mathbf{x}} = {} & \lambda_r v_r + \lambda_\theta\frac{v_\theta}{r} + \lambda_{v_r}\left(\frac{v_\theta^2}{r} - \frac{\mu}{r^2} + a_r\right) \nonumber\\
& + \lambda_{v_\theta}\left(-\frac{v_r v_\theta}{r} + a_\theta\right) - \lambda_m\frac{m^2}{2P}(a_r^2+a_\theta^2).
\end{align}
Adding the Lagrangian $L=1$ gives the complete Hamiltonian for the sun-field model. For a minimum-time problem, the PMP asserts that along an optimal trajectory:
\begin{enumerate}
    \item The state and costates evolve according to Hamilton's
          canonical equations: $\dot{\mathbf{x}} = \partial\mathcal{H}/\partial\boldsymbol{\lambda}$
          and $\dot{\boldsymbol{\lambda}} = -\partial\mathcal{H}/\partial\mathbf{x}$.
    \item The optimal control $\mathbf{a}_T(t)$ maximizes $\mathcal{H}$ point-wise.
    \item The Hamiltonian remains constant and equal to zero: $\mathcal H\equiv 0$.
\end{enumerate}

\subsection{Optimal Control Law}
To find the control that maximizes $\mathcal{H}$, we set the partial derivatives of the Hamiltonian with respect to the control variables $(a_r, a_\theta)$ to zero:
\begin{subequations}
\begin{align}
    \frac{\partial\mathcal{H}}{\partial a_r} &= \lambda_{v_r} - \lambda_m \frac{m^2}{P} a_r = 0, \\
    \frac{\partial\mathcal{H}}{\partial a_\theta} &= \lambda_{v_\theta} - \lambda_m \frac{m^2}{P} a_\theta = 0.
\end{align}
\end{subequations}
This yields the optimal thrust acceleration vector:
\begin{equation}
  {\bf a}_T = k(t)\,\boldsymbol\lambda_v, \quad \text{where} \quad \boldsymbol\lambda_v \equiv (\lambda_{v_r},\lambda_{v_\theta})^\top,
  \label{eq:ThrustLaw_prelim}
\end{equation}
and the time-varying gain factor is defined as
\begin{equation}
    k(t) \equiv \frac{P}{\lambda_m(t) m^2(t)}.
    \label{eq:GainFactorDef}
\end{equation}
The optimal thrust direction is parallel to the velocity costate vector $\boldsymbol\lambda_v$. We now show that the gain factor $k(t)$ is, in fact, a constant of motion.
Recent adaptive algorithms for low-thrust trajectory design exploit similar primer-vector structures to generate high-quality initial guesses for complex transfers \cite{PinoHowell2020}.

\subsection{Costate Equations and First Integrals}
The costate dynamics are found from $\dot\lambda_i = -\partial\mathcal{H}/\partial x_i$:
\begin{subequations}\label{eq:CostateDyn}
\begin{align}
  \dot\lambda_r &=
      \frac{\lambda_\theta v_\theta}{r^{2}}
      +\frac{\lambda_{v_r} v_\theta^{2}}{r^{2}}
      -2\lambda_{v_r}\frac{\mu}{r^{3}}
      -\lambda_{v_\theta}\frac{v_r v_\theta}{r^{2}},
      \label{eq:lam_r}\\[4pt]
  \dot\lambda_\theta &= -\frac{\partial\mathcal{H}}{\partial\theta} = 0,
      \label{eq:lam_theta_deriv}\\[4pt]
  \dot\lambda_{v_r} &=
      -\lambda_r + \lambda_{v_\theta}\frac{v_\theta}{r},
      \label{eq:lam_vr}\\[4pt]
  \dot\lambda_{v_\theta} &=
      -\frac{\lambda_\theta}{r}
      -2\lambda_{v_r}\frac{v_\theta}{r}
      +\lambda_{v_\theta}\frac{v_r}{r},
      \label{eq:lam_vt}\\[4pt]
  \dot\lambda_m &= -\frac{\partial\mathcal{H}}{\partial m} =
      \lambda_m \frac{m}{P}\bigl(a_r^{2}+a_\theta^{2}\bigr).
      \label{eq:lam_m}
\end{align}
\end{subequations}
From these dynamics, we extract two first integrals.

\paragraph{Rotational Symmetry}
Equation \eqref{eq:lam_theta_deriv} arises because the Hamiltonian is independent of the absolute angle $\theta$ (a cyclic coordinate). This gives our first integral:
\begin{equation}
  \boxed{\lambda_\theta(t) = C_\theta = \text{const}}.
\end{equation}

\paragraph{Mass–Costate Relationship}
We find a direct relationship between mass $m$ and its costate $\lambda_m$ by examining their rates of change. Using the chain rule, we have $\frac{{\rm d}\lambda_m}{{\rm d}m} = \frac{\dot\lambda_m}{\dot m}$. Substituting from \eqref{eq:lam_m} and \eqref{eq:StateDyn}e:
\begin{equation}
    \frac{{\rm d}\lambda_m}{{\rm d}m} = \frac{\lambda_m \frac{m}{P}\bigl(a_r^{2}+a_\theta^{2}\bigr)}{-\frac{m^{2}}{2P}\bigl(a_r^{2}+a_\theta^{2}\bigr)} = \frac{\lambda_m m}{-m^2/2} = -2\frac{\lambda_m}{m}. \nonumber
\end{equation}
This is a separable first-order ordinary differential equation:
\[
    \frac{{\rm d}\lambda_m}{\lambda_m} = -2\frac{{\rm d}m}{m}.
\]
Integrating both sides yields $\ln(\lambda_m) = -2\ln(m) + \text{const}$, which upon exponentiation gives our second, crucial integral of motion:
\begin{equation}
  \boxed{\lambda_m(t) m^2(t) = C = \text{const}}.
  \label{eq:IntMass}
\end{equation}
This relation holds for the entire trajectory.

\paragraph{Constant Thrust Gain}
Substituting the integral from \eqref{eq:IntMass} back into the definition of the gain factor \eqref{eq:GainFactorDef}, we find that it is constant:
\begin{equation}
  k(t) = \frac{P}{\lambda_m(t) m^2(t)} = \frac{P}{C} \equiv k_T = \text{const}.
\end{equation}
The optimal control law \eqref{eq:ThrustLaw_prelim} therefore simplifies to:
\begin{equation}
  \mathbf{a}_T(t) = k_T\,\boldsymbol{\lambda}_v(t).
  \label{eq:ThrustLaw_Final}
\end{equation}
The thrust magnitude is proportional to the norm of the velocity costates, $\|\boldsymbol{\lambda}_v\|$, and the proportionality factor $k_T$ is a constant determined by the boundary conditions. The discovery of this constant gain is a key analytical result, as it significantly simplifies the nature of the optimal control.

\medskip
This completes the analytical formulation. The problem is reduced to a boundary-value problem governed by the seven differential equations for the state $(r, \theta, v_r, v_\theta)$ and non-constant costates $(\lambda_r, \lambda_{v_r}, \lambda_{v_\theta})$, with the control given by \eqref{eq:ThrustLaw_Final}. A shooting method must find the initial values of the four costates $(\lambda_r(0), C_\theta, \lambda_{v_r}(0), \lambda_{v_\theta}(0))$ and the thrust constant $k_T$ that satisfy the specified final conditions at the optimal final time $t_f$.

\section{Similarity Reduction of the Sun-Field Transfer Family}
\label{sec:similarity_reduction}

The Sun-field boundary-value problem from Sec.~\ref{sec:SunModel} is written in dimensional variables, but its family of circular-to-circular transfers possesses a useful scaling symmetry. 
This symmetry allows the six nominal inputs
\[
(r_0,\; r_f,\; \Theta,\; P,\; m_i,\; m_\mathrm{dry})
\]
to be reduced to three dimensionless parameters. 
The reduction is useful both conceptually and numerically: conceptually, because it shows which combinations of the physical inputs actually determine the transfer shape; numerically, because it makes it possible to tabulate a reduced atlas of normalized solutions and then map them back to dimensional trajectories for any specific spacecraft.

\subsection{Canonical units and normalized variables}
Let the departure radius define the distance unit,
\begin{equation}
  R_0 \equiv r_0,
\end{equation}
and let the associated Keplerian time and velocity scales be
\begin{equation}
  T_0 \equiv \sqrt{\frac{R_0^3}{\mu}},
  \qquad
  V_0 \equiv \frac{R_0}{T_0} = \sqrt{\frac{\mu}{R_0}}.
\end{equation}
The corresponding acceleration scale is
\begin{equation}
  A_0 \equiv \frac{V_0}{T_0} = \frac{\mu}{R_0^2}.
\end{equation}

We now define the dimensionless independent variable and normalized state by
\begin{equation}
  \tau \equiv \frac{t}{T_0},
  \qquad
  \bar r \equiv \frac{r}{R_0},
  \qquad
  \bar v_r \equiv \frac{v_r}{V_0},
  \qquad
  \bar v_\theta \equiv \frac{v_\theta}{V_0},
\end{equation}
and similarly normalize the thrust acceleration components as
\begin{equation}
  \bar a_r \equiv \frac{a_r}{A_0},
  \qquad
  \bar a_\theta \equiv \frac{a_\theta}{A_0}.
\end{equation}
The polar angle \(\theta\) is already dimensionless and is left unchanged.

The mass is handled more conveniently through the inverse-mass increment that already appears in the free-space solution. 
Define
\begin{equation}
  \Delta \equiv \frac{1}{m_\mathrm{dry}}-\frac{1}{m_i},
\end{equation}
and introduce the normalized fuel coordinate
\begin{equation}
  q(\tau) \equiv \frac{m(\tau)^{-1}-m_i^{-1}}{\Delta}.
  \label{eq:q_def_similarity}
\end{equation}
By construction,
\begin{equation}
  q(0)=0,
  \qquad
  q(\tau_f)=1,
\end{equation}
so \(q\) measures the fraction of the available inverse-mass budget that has been consumed. 
The dimensional mass is recovered from \(q\) as
\begin{equation}
  m(\tau)=\frac{1}{m_i^{-1}+\Delta q(\tau)}.
  \label{eq:m_from_q_similarity}
\end{equation}

\subsection{Dimensionless form of the dynamics}
Substituting the scaled variables into \eqref{eq:StateDyn} gives
\begin{subequations}\label{eq:StateDyn_scaled}
\begin{align}
  \bar r' &= \bar v_r, \\
  \theta' &= \frac{\bar v_\theta}{\bar r}, \\
  \bar v_r' &= \frac{\bar v_\theta^2}{\bar r} - \frac{1}{\bar r^2} + \bar a_r, \\
  \bar v_\theta' &= -\frac{\bar v_r\bar v_\theta}{\bar r} + \bar a_\theta, \\
  q' &= \frac{\bar a_r^2+\bar a_\theta^2}{\kappa},
\end{align}
\end{subequations}
where prime denotes differentiation with respect to \(\tau\), and where the single dimensionless propulsion parameter is
\begin{equation}
  \kappa \equiv 2P\Delta\,\frac{R_0^{5/2}}{\mu^{3/2}}
  = 2P\left(\frac{1}{m_\mathrm{dry}}-\frac{1}{m_i}\right)\frac{r_0^{5/2}}{\mu^{3/2}}.
  \label{eq:kappa_def}
\end{equation}

This is the key reduction. 
In dimensional form, the mass-flow law \eqref{eq:mass_flow} depends separately on \(P\), \(m_i\), \(m_\mathrm{dry}\), and \(r_0\). 
After scaling, those four quantities appear only through the single combination \(\kappa\). 
The normalized equations therefore no longer distinguish between spacecraft that have different powers and masses but the same value of \(\kappa\). 
Such spacecraft follow the same dimensionless trajectory and differ only by the dimensional scales used to reconstruct time, radius, velocity, and acceleration.

\subsection{Dimensionless boundary conditions}
For a transfer from a circular orbit of radius \(r_0\) to a circular orbit of radius \(r_f\), the normalized boundary conditions become
\begin{subequations}\label{eq:BC_scaled}
\begin{align}
  \bar r(0) &= 1, &
  \bar r(\tau_f) &= \rho, \\
  \theta(0) &= 0, &
  \theta(\tau_f) &= \Theta,\\
  \bar v_r(0) &= 0, &  \bar v_r(\tau_f) &= 0, \\
  \bar v_\theta(0) &= 1, &\bar v_\theta(\tau_f) &= \rho^{-1/2}, \\
  q(0) &= 0, &q(\tau_f) &= 1,  
\end{align}
\end{subequations}
with
\begin{equation}
  \rho \equiv \frac{r_f}{r_0}.
  \label{eq:rho_def_similarity}
\end{equation}
The quantity \(\rho\) is the radius ratio, while
\begin{equation}
  \Theta \equiv \theta_f-\theta_0+2\pi N
  \label{eq:Theta_def_similarity}
\end{equation}
is the unwrapped transfer angle, with \(N\in\mathbb Z_{\ge 0}\) selecting the number of completed revolutions. 
Only the angular difference matters, not the absolute departure angle, because the Sun-field model is rotationally invariant. 
This is the same symmetry that produced the first integral \(\lambda_\theta=C_\theta\) in Sec.~\ref{sec:SunModel}.

Equations~\eqref{eq:StateDyn_scaled}--\eqref{eq:BC_scaled} show that the normalized transfer problem depends only on
\[
(\rho,\;\kappa,\;\Theta).
\]
This is the reduced parameter space.

\subsection{Why the parameter space shrinks}
It is useful to state explicitly what has been removed.

First, the absolute orbital scale \(r_0\) disappears from the normalized equations because lengths are measured in units of \(r_0\), times in units of \(\sqrt{r_0^3/\mu}\), and velocities in units of \(\sqrt{\mu/r_0}\). 
Only the ratio \(r_f/r_0\) remains.

Second, the three propulsion quantities \(P\), \(m_i\), and \(m_\mathrm{dry}\) do not independently affect the normalized dynamics. 
They enter only through the inverse-mass budget \(\Delta\) and then only through the combination \(\kappa\) in \eqref{eq:kappa_def}. 
Thus three nominal propulsion parameters collapse to one similarity parameter.

Third, the absolute phase angle \(\theta_0\) is irrelevant because the equations are invariant under a rigid rotation of the entire trajectory. 
Only the total angular change \(\Theta\) matters.

As a result, the circular-to-circular Sun-field transfer problem is parameterized, at the similarity level, by the three variables \((\rho,\kappa,\Theta)\) rather than the six nominal quantities \((r_0,r_f,\Theta,P,m_i,m_\mathrm{dry})\).


\subsection{Forward map from nominal parameters to similarity space}
For a given dimensional transfer request, the reduced coordinates are obtained directly from the nominal mission parameters:
\begin{align}
  \rho &= \frac{r_f}{r_0}, \\
  \Theta &= \theta_f-\theta_0+2\pi N, \\
  \Delta &= \frac{1}{m_\mathrm{dry}}-\frac{1}{m_i}, \\
  \kappa &= 2P\Delta\,\frac{r_0^{5/2}}{\mu^{3/2}}.
\end{align}
Here \(N\) selects the desired multi-revolution branch. 
Thus the similarity reduction is applied by first computing \((\rho,\kappa,\Theta)\) from the physical problem and then solving the corresponding normalized boundary-value problem.


\subsection{Inverse map from a normalized atlas solution to dimensional variables}
Suppose now that a normalized solution has been computed or interpolated in the reduced space \((\rho,\kappa,\Theta)\). 
That solution consists of the dimensionless functions
\[
\bar r(\tau),\quad
\theta(\tau),\quad
\bar v_r(\tau),\quad
\bar v_\theta(\tau),\quad
q(\tau),\quad
\bar a_r(\tau),\quad
\bar a_\theta(\tau),
\]
together with the normalized final time \(\tau_f\).

For the specific dimensional system defined by \(r_0\), \(m_i\), \(m_\mathrm{dry}\), and \(\mu\), the trajectory is recovered by the inverse scaling
\begin{align}
  t &= T_0\tau = \sqrt{\frac{r_0^3}{\mu}}\,\tau, \\
  r(t) &= r_0\,\bar r(\tau), \\
  \theta(t) &= \theta_0 + \theta(\tau), \\
  v_r(t) &= \sqrt{\frac{\mu}{r_0}}\,\bar v_r(\tau), \\
  v_\theta(t) &= \sqrt{\frac{\mu}{r_0}}\,\bar v_\theta(\tau), \\
  a_r(t) &= \frac{\mu}{r_0^2}\,\bar a_r(\tau), \\
  a_\theta(t) &= \frac{\mu}{r_0^2}\,\bar a_\theta(\tau), \\
  m(t) &= \frac{1}{m_i^{-1}+\Delta q(\tau)}.
\end{align}
In particular, the dimensional flight time is
\begin{equation}
  t_f = \sqrt{\frac{r_0^3}{\mu}}\,\tau_f.
\end{equation}

This gives the practical recipe for using a precomputed solution atlas:
\begin{enumerate}
    \item Start from the nominal physical parameters.
    \item Compute \(\rho\), \(\Theta\), and \(\kappa\) from the forward map above.
    \item Retrieve or interpolate the normalized solution stored at that location in \((\rho,\kappa,\Theta)\)-space.
    \item Recover the dimensional trajectory by the inverse scaling formulas.
\end{enumerate}

\subsection{What should be stored in the atlas}
For numerical continuation, it is not necessary to store the full dimensional trajectory. 
A more efficient choice is to store the normalized shooting data, namely the initial costates, the constant \(C_\theta\) associated with rotational symmetry, the thrust-gain constant \(k_T\), and the normalized flight time. 
Because these quantities are stored in dimensionless form, the same atlas entry can be reused for any dimensional spacecraft that maps to the same \((\rho,\kappa,\Theta)\). 
This is the computational advantage of the similarity reduction: one normalized solution represents an entire equivalence class of physical missions.


\section{Numerical Procedure}
Solving the two-point boundary-value problem requires a numerical approach. We employ a two-stage strategy that combines a global search to find an initial, feasible trajectory with a local refinement process that uses a continuation method to precisely reach the desired target \cite{Fogel2020,DiPasquale2024}.

\subsection{Objective Function}
The core of the numerical method is an objective function, $J$, which quantifies the quality of a given trajectory. This function is minimized to find the optimal set of initial costate parameters. It incorporates penalties for violating physical constraints and errors for missing the target state. The function has two distinct forms depending on the optimization stage.

\subsubsection{Global Search Objective}
In the initial search phase, the goal is to find \emph{any} valid transfer to a stable circular orbit, without a specific target location. The objective function, $J_\text{global}$, is formulated to reward reaching a circular velocity in the outer solar system while penalizing physically unrealistic outcomes:
\begin{equation}
    J_\text{global} = w_f \mathcal{P}_f + w_r \mathcal{P}_r + w_v \mathcal{E}_v^2.
    \label{eq:obj_global}
\end{equation}
The terms are defined as:
\begin{itemize}
    \item $\mathcal{P}_f$: A fuel penalty that applies a large constant value if the final mass $m_f$ drops below the dry mass $m_\text{dry}$, effectively creating a hard constraint against using more fuel than available.
    \item $\mathcal{P}_r$: A radial position penalty, $\max(0, r_\text{min} - r_f)^2$, which discourages solutions that fail to travel a minimum distance $r_\text{min}$ (e.g., 1.5\,AU) from the Sun.
    \item $\mathcal{E}_v$: The velocity mismatch error at the final state. It measures the Euclidean distance between the final velocity vector $\mathbf{v}_f=(v_{r,f}, v_{\theta,f})$ and the required velocity for a circular orbit at the final radius $r_f$, given by $\mathbf{v}_\text{circ}=(0, \sqrt{\mu/r_f})$. The squared error is $\mathcal{E}_v^2 = v_{r,f}^2 + (v_{\theta,f} - \sqrt{\mu/r_f})^2$.
    \item $w_f, w_r, w_v$: Positive weighting constants that balance the contribution of each term.
\end{itemize}

\subsubsection{Local Refinement Objective}
Once an initial solution is found, the objective function is switched to a targeted version, $J_\text{local}$, to steer the trajectory to a precise destination $(r_\text{des}, \theta_\text{des})$:\
\begin{equation}
    J_\text{local} = w_f \mathcal{P}_f + w_v \mathcal{E}_v^2 + w_r (r_f - r_\text{des})^2 + w_\theta (\theta_f - \theta_\text{des})^2.
    \label{eq:obj_local}
\end{equation}
Here, the radial penalty $\mathcal{P}_r$ is removed, and two new terms are added to penalize deviations from the desired final radius $r_\text{des}$ and true anomaly $\theta_\text{des}$.

\subsection{Two-Stage Optimization Strategy}
\subsubsection{Stage 1: Global Search}
An initial, admissible solution is found by minimizing the global objective function \eqref{eq:obj_global}. For this task, we use the \texttt{differential\_evolution} algorithm from the SciPy library \cite{Storn1997}. This robust, population-based method is well-suited for exploring the wide parameter space of the initial costates without getting trapped in local minima, thereby finding a good starting point for the next stage.

\subsubsection{Stage 2: Continuation and Local Refinement}
Starting from the solution found in Stage~1, which ends at a point $(r^\star, \theta^\star)$, we use a continuation (or homotopy) method to incrementally approach the final desired target $(r_\text{des}, \theta_\text{des})$ \cite{Allgower1990}. The procedure is as follows:
\begin{enumerate}
    \item Set an intermediate target partway between the current known solution's endpoint and the final desired endpoint.
    \item Perform a local optimization by minimizing the targeted objective function \eqref{eq:obj_local} using Powell's method \cite{Powell1964} (derivative-free) as implemented in \texttt{scipy.optimize.minimize}. The initial guess for the optimizer is the set of parameters from the previous successfully converged solution.
    \item If the optimizer converges to a high-precision solution, the new solution is accepted. The intermediate target is then advanced closer to the final destination, and the process repeats.
    \item If the optimizer fails to converge, it indicates the step to the new target was too large. The intermediate target is then moved back, typically halving the distance between the last successful endpoint and the current target, and the optimization is re-attempted.
\end{enumerate}
This adaptive stepping ensures that the local optimizer is always given a problem that is a small perturbation from a known good solution, guaranteeing robust convergence across the entire solution space until the final target is reached.

All simulations, optimization procedures, and figures were generated using the accompanying open-source Python implementation \cite{Olsina2025}.
The repository includes the complete solver, plotting scripts, and resulting graphs of trajectories, fuel and acceleration magnitudes.

\pagebreak

\section{Results}
The numerical procedure was applied to compute minimum-time transfers from a 1\,AU circular orbit (representing Earth) to various target circular orbits. The targets are defined by their final radius and true anomaly relative to the starting position. A 1\,GW engine was assumed with an initial vehicle mass of 3000\,tons and a dry mass of 1000\,tons.

Figures~\ref{fig:final_theta} and~\ref{fig:starting_theta} show the relationship between flight time and the orbital geometry. The flight times exhibit a U-shaped dependence on the final true anomaly, with a clear minimum for transfers that are neither direct short-arcs nor long-arcs. Similar behavior has been reported in recent multi-objective studies of variable-$I_{\!sp}$ maneuvers \cite{DiPasquale2024}.

Figures~\ref{fig:oberth_figure} and~\ref{fig:oberth_figure2} show example trajectories. A key feature of these optimal paths is the inclusion of an Oberth maneuver. To gain orbital energy efficiently, the spacecraft first thrusts to lower its periapsis, making a close pass of the Sun. It then applies its main acceleration burn at maximum velocity near the Sun, which is far more effective at increasing orbital energy than thrusting in a higher orbit \cite{Oberth1923}. This is visible in the plots as an initial dip towards the Sun and a corresponding peak in the acceleration magnitude.

Table~\ref{tab:transfer_times_final_anomaly} and Table~\ref{tab:transfer_times_initial_anomaly} present the resulting transfer times in days for a grid of targets. As expected, transfers to more distant orbits take longer. The notation 'N/A' indicates that the transfer was not achievable with the given propellant load. This typically happens because there is a fixed constraint that cannot be overcome by continuous change of the initial conditions. Figure \ref{fig:oberth_figure2} shows trajectory  very close to such a barrier, with very extreme Oberth maneuver around the Sun that cannot be extended further. 


\section{Conclusion}
We combined analytic insight (two closed-form first integrals and a free-space benchmark solution) with a robust continuation strategy to solve the time-optimal transfer problem for a constant-power, variable-$I_{\!sp}$ engine. The method successfully generates high-performance trajectories, including those featuring Oberth maneuvers, for transfers out to the giant planets with a gigawatt-class engine. Recent developments in variable-$I_{\!sp}$ propulsion and adaptive low-thrust guidance algorithms \cite{PinoHowell2020,DiPasquale2024} suggest that further improvements in mission performance are possible. Future work will extend this framework to three-dimensional rendezvous problems.

\section*{Acknowledgment}
The author wishes to acknowledge the use of large language models (Google's Gemini and OpenAI's GPT) for assistance in improving the language and clarity of the manuscript, for aiding in the implementation of the numerical code and derivations of some of the analytical results. The results were tested and verified by the author.

\begin{figure}[htbp]
    \centering
    \includegraphics[width=\linewidth]{article_003.jpg}
    \caption{Flight times in days versus final true anomaly at arrival.}
    \label{fig:final_theta}
\end{figure}
\begin{figure}[htbp]
    \centering
    \includegraphics[width=\linewidth]{article_004.jpg}
    \caption{Flight times in days versus initial true anomaly at departure.}
    \label{fig:starting_theta}
\end{figure}

\begin{figure*}[p]
    \centering
    \includegraphics[width=0.75\linewidth]{article_001.jpg}
    \caption{Example trajectory from $r=1$\,AU, $\theta=0^\circ$ to $r=5$\,AU, $\theta=273.1^\circ$ in 207 days. The trajectory includes a solar flyby (Oberth maneuver) to gain energy, which is apparent from the acceleration peak around day 50.}
    \label{fig:oberth_figure}
\end{figure*}
\begin{figure*}[p]
    \centering
    \includegraphics[width=\linewidth]{article_005.jpg}
    \caption{Example trajectory from $r=1$\,AU, $\theta=0^\circ$ to $r=1.5$\,AU, $\theta=34.7^\circ$ in 44 days. The trajectory is a very close analogy to the trajectories without gravity, demonstrating clear acceleration-deceleration pattern that is almost fully symmetric.}
    \label{fig:non_oberth_figure}
\end{figure*}
\begin{figure*}[p]
    \centering
    \includegraphics[width=\textwidth]{article_002.jpg}
    \caption{An example of a more extreme Oberth maneuver for a transfer to Jupiter's orbit. The trajectory is from $r=1$\,AU, $\theta=0^\circ$ to $r=5.2$\,AU, $\theta=385.2^\circ$ and lasts 214 days.}
    \label{fig:oberth_figure2}
\end{figure*}


\FloatBarrier


\appendices
\section{Power-System Mass Closure for the Free-Space Benchmark}
\label{sec:appendix_mass_closure}

The benchmark result \eqref{eq:Tmin} was derived for prescribed initial and dry masses. 
For high-power spacecraft, however, the dry mass may itself depend on the propulsion system, because both the power-conversion hardware and the radiators must be carried together with the payload. 
This appendix shows how the free-space minimum-time law is modified under a simple power-system mass closure.

Let \(P\) denote the useful jet power appearing in the rocket model, and let \(P_\mathrm{tot}\) be the total onboard engine power. 
Define
\begin{equation}
  \alpha_\mathrm{eng} \equiv \frac{P_\mathrm{tot}}{m_\mathrm{eng}},
  \qquad
  \phi \equiv \frac{P_\mathrm{heat}}{P_\mathrm{tot}},
  \qquad
  \rho_\mathrm{rad} \equiv \frac{P_\mathrm{heat}}{m_\mathrm{rad}},
\end{equation}
where \(\alpha_\mathrm{eng}\) is the engine specific power, \(\phi\) is the waste-heat fraction, and \(\rho_\mathrm{rad}\) is the radiator specific heat-rejection capability. 
Since only the fraction \(1-\phi\) of the total power is converted into useful jet power,
\begin{equation}
  P = (1-\phi)P_\mathrm{tot},
  \qquad
  P_\mathrm{heat} = \phi P_\mathrm{tot}
  = \frac{\phi}{1-\phi}P.
\end{equation}
Therefore,
\begin{equation}
  m_\mathrm{eng}=\frac{P}{(1-\phi)\alpha_\mathrm{eng}},
  \qquad
  m_\mathrm{rad}=\frac{\phi}{1-\phi}\frac{P}{\rho_\mathrm{rad}}.
\end{equation}

Let \(m_\mathrm{pay}\) be the payload mass. 
The dry mass can then be written in the affine form
\begin{equation}
  m_\mathrm{dry}=m_\mathrm{pay}+m_\mathrm{eng}+m_\mathrm{rad}
  = m_\mathrm{pay}+\beta P,
  \label{eq:mdry_beta_app}
\end{equation}
with
\begin{equation}
  \beta \equiv
  \frac{1}{(1-\phi)\alpha_\mathrm{eng}}
  +\frac{\phi}{(1-\phi)\rho_\mathrm{rad}}.
  \label{eq:beta_def_app}
\end{equation}
The quantity \(\beta\) has units of kg/W and is the inverse of the effective propulsion-system specific power,
\begin{equation}
  \alpha_\mathrm{sys}\equiv \beta^{-1}.
\end{equation}

To close the mass model, assume a fixed fuel-to-dry-mass ratio
\begin{equation}
  \mathcal{R}\equiv \frac{m_\mathrm{fuel}}{m_\mathrm{dry}}.
\end{equation}
Then
\begin{equation}
  m_i=(1+\mathcal{R})m_\mathrm{dry},
\end{equation}
and the parameter \(\Delta\) from \eqref{eq:Tmin} becomes
\begin{equation}
  \Delta
  = \frac{1}{m_\mathrm{dry}}-\frac{1}{m_i}
  = \frac{\mathcal{R}}{(1+\mathcal{R})m_\mathrm{dry}}.
  \label{eq:Delta_payload_app}
\end{equation}
Substituting \eqref{eq:mdry_beta_app} and \eqref{eq:Delta_payload_app} into the benchmark law \eqref{eq:Tmin} yields
\begin{equation}
  T(P)
  =
  \left[
    \frac{6L^2(1+\mathcal{R})}{\mathcal{R}}
    \left(
      \frac{m_\mathrm{pay}}{P}+\beta
    \right)
  \right]^{1/3}.
  \label{eq:T_payload_beta_app}
\end{equation}
Equation~\eqref{eq:T_payload_beta_app} is the free-space minimum-time law when the dry mass is allowed to scale linearly with propulsion power.

Several immediate consequences follow. 
First,
\begin{equation}
  T^3(P)
  =
  \frac{6L^2(1+\mathcal{R})}{\mathcal{R}}
  \left(
    \frac{m_\mathrm{pay}}{P}+\beta
  \right),
\end{equation}
so that
\begin{equation}
  \frac{d}{dP}T^3(P)
  =
  -\frac{6L^2(1+\mathcal{R})}{\mathcal{R}}
   \frac{m_\mathrm{pay}}{P^2}
  <0.
\end{equation}
Under the idealized assumption of unlimited exhaust speed, the transfer time therefore decreases monotonically with power and has no finite mathematical optimum:
\begin{equation}
  P^\star \to \infty.
\end{equation}
Second, the family approaches the asymptotic lower bound
\begin{equation}
  T_\infty
  =
  \left[
    \frac{6L^2(1+\mathcal{R})}{\mathcal{R}}\,\beta
  \right]^{1/3}.
  \label{eq:T_infty_beta_app}
\end{equation}
This lower bound arises because increasing \(P\) also increases the dry mass that must be accelerated.

A convenient engineering scale is obtained by equating propulsion-hardware mass to payload mass,
\begin{equation}
  m_\mathrm{eng}+m_\mathrm{rad}=m_\mathrm{pay}.
\end{equation}
Using \eqref{eq:mdry_beta_app}, this defines the characteristic power
\begin{equation}
  \beta P_\mathrm{knee}=m_\mathrm{pay},
  \qquad
  P_\mathrm{knee}=\frac{m_\mathrm{pay}}{\beta}
  = \alpha_\mathrm{sys}m_\mathrm{pay}.
  \label{eq:Pknee_app}
\end{equation}
At this point the dry mass is \(m_\mathrm{dry}=2m_\mathrm{pay}\), and \eqref{eq:T_payload_beta_app} can be written as
\begin{equation}
  T(P)=T_\infty
  \left(
    1+\frac{P_\mathrm{knee}}{P}
  \right)^{1/3}.
  \label{eq:T_knee_form_app}
\end{equation}
Hence
\begin{equation}
  T(P_\mathrm{knee}) = 2^{1/3}T_\infty.
  \label{eq:Tknee_app}
\end{equation}

Equations~\eqref{eq:T_payload_beta_app}--\eqref{eq:Tknee_app} provide an engineering interpretation of the free-space benchmark. 
They are not used in the subsequent Sun-field optimal-control formulation, but they show how the cubic-root time law is modified when propulsion-system mass scales with power.

\section{Direct Method Using Spline Parameterization}

\label{sec:appendix}
Before adopting the indirect method based on Pontryagin's Maximum Principle (PMP) detailed in the main body of this paper, a more straightforward \emph{direct method} was explored. Direct methods seek to find the optimal trajectory by parameterizing the path with a set of functions and then using a nonlinear programming solver to find the best parameters. This approach avoids the complexity of deriving and solving the costate equations.

Our implementation parameterized the spacecraft's trajectory, $\mathbf{r}(t) = (x(t), y(t))$, using cubic splines. The core idea is to define the path by a small number of control points, which become the variables for optimization.

The procedure was as follows:
{
\renewcommand{\labelenumi}{\arabic{enumi}.}
\begin{enumerate}
    \item A fixed number of $N$ control points, $(\mathbf{p}_i, t_i)$ for $i=0, \dots, N-1$, were defined along the time axis of the transfer. The endpoints, $\mathbf{p}_0$ and $\mathbf{p}_{N-1}$, were fixed to match the start and end positions of the transfer arc. The initial and final velocities were enforced by connecting the transfer arc to the initial and final circular orbits.
    \item The positions of the intermediate $N-2$ control points were the free variables for the optimization.
    \item For any given set of control points, two cubic splines, $x_s(t)$ and $y_s(t)$, were fitted to create a continuous and smooth trajectory $\mathbf{r}(t) = (x_s(t), y_s(t))$. A cubic spline is defined piecewise by third-order polynomials, ensuring that the first and second derivatives (velocity and acceleration) are continuous across the entire path.
    \item The required velocity $\dot{\mathbf{r}}(t)$ and acceleration $\ddot{\mathbf{r}}(t)$ vectors were obtained by numerically differentiating the spline functions.
    \item With the total acceleration $\ddot{\mathbf{r}}(t)$ known, the required thrust acceleration was calculated by subtracting the gravitational component:
    \begin{equation}
        \mathbf{a}_T(t) = \ddot{\mathbf{r}}(t) - \mathbf{a}_g(\mathbf{r}(t)).
    \end{equation}
    \item The instantaneous mass flow rate was then determined from the thrust acceleration magnitude using \eqref{eq:mass_flow}. Integrating this rate backwards in time yielded the total propellant consumed for the maneuver.
    \item A nonlinear optimizer (specifically, the Nelder--Mead method) was used to adjust the positions of the intermediate control points to minimize a cost function based on the final fuel expenditure and flight time.
\end{enumerate}
}

\subsection*{Limitations and Stagnation in Local Minima}
While conceptually simple, this direct method proved inadequate for finding high-performance trajectories. The primary difficulty was the nature of the optimization landscape. The relationship between the position of a control point and the final mission cost is highly complex and nonlinear. As a result, the search space is filled with a vast number of local minima.

For a very small number of intermediate control points ($N=3$ or $N=4$ total), the method could find a reasonable, smooth trajectory. However, as the number of control points was increased to grant the path more freedom and potential for higher performance, the dimensionality of the search space grew, and the solver would invariably become trapped in a suboptimal local minimum.

An attempt was made to mitigate this by using a continuation strategy on the number of control points: first, an optimal path with $N=4$ points was found. Then, a fifth point was added at the location of maximum curvature, and the system was re-optimized. While this yielded marginal improvements, the solution quality consistently stagnated far below that achievable with the PMP-based method. The direct spline method lacks the rigorous guidance provided by the costate (or "primer vector") variables, which point the way toward optimality across the entire trajectory. These difficulties ultimately motivated the adoption of the indirect shooting method presented in this work.

\bibliographystyle{ieeetr}

\bibliography{references}

%\begin{table*}[!htbp]
%\centering
%\caption{Minimum transfer time (in days) to a circular orbit of a given %radius and final true anomaly.}
%\label{tab:transfer_times}
%\sisetup{table-format=3.2}
%\begin{tabular}{r S S S S S S S S}
%\toprule
%\textbf{True Anomaly} & \multicolumn{8}{c}{\textbf{Target Radius (AU)}} %\\
%\cmidrule(lr){2-9}
%\textbf{(\si{\degree})} & {0.3} & {0.75} & {1.643} & {3} & {5} & {10} & %{15} & {20} \\
%\midrule
%-100 & {N/A} & {N/A} & 172.79 & 214.70 & 278.24 & 411.05 & 521.03 & 617.79 \\
%-95  & {N/A} & {N/A} & 172.53 & 214.38 & 278.09 & 411.36 & 521.71 & 618.81 \\
%-90  & {N/A} & {N/A} & 172.04 & 213.82 & 277.69 & 411.37 & 522.08 & 619.49 \\
%-85  & {N/A} & {N/A} & 171.31 & 213.01 & 277.01 & 411.07 & 522.09 & 619.81 \\
%-80  & {N/A} & {N/A} & 170.30 & 211.89 & 276.00 & 410.38 & 521.69 & 619.66 \\
%-75  & {N/A} & {N/A} & 169.02 & 210.48 & 274.68 & 409.32 & 520.91 & 619.10 \\
%-70  & {N/A} & {N/A} & 167.45 & 208.77 & 273.04 & 407.90 & 519.68 & 618.09 \\
%-65  & {N/A} & {N/A} & 165.56 & 206.75 & 271.06 & 406.11 & 518.00 & 616.62 \\
%-60  & {N/A} & {N/A} & 163.31 & 204.37 & 268.70 & 403.85 & 515.90 & 614.62 \\
%-55  & {N/A} & {N/A} & 160.72 & 201.64 & 265.98 & 401.20 & 513.34 & 612.07 \\
%-50  & {N/A} & {N/A} & 157.75 & 198.53 & 262.91 & 398.15 & 510.26 & 609.06 \\
%-45  & {N/A} & {N/A} & 154.36 & 195.05 & 259.48 & 394.69 & 506.73 & 605.55 \\
%-40  & {N/A} & {N/A} & 150.51 & 191.19 & 255.65 & 390.78 & 502.75 & 601.49 \\
%-35  & {N/A} & {N/A} & 146.19 & 186.97 & 251.46 & 386.49 & 498.35 & 596.95 \\
%-30  & {N/A} & {N/A} & 141.39 & 182.31 & 246.95 & 381.83 & 493.47 & 591.90 \\
%-25  & {N/A} & {N/A} & 136.01 & 177.29 & 242.13 & 376.83 & 488.23 & 586.47 \\
%-20  & {N/A} & {N/A} & 129.99 & 171.93 & 236.97 & 371.49 & 482.60 & 580.58 \\
%-15  & {N/A} & {N/A} & 123.33 & 166.21 & 231.57 & 365.86 & 476.66 & 574.36 \\
%-10  & {N/A} & {N/A} & 115.97 & 160.22 & 225.98 & 360.03 & 470.45 & 567.80 \\
%-5   & {N/A} & {N/A} & 107.80 & 154.01 & 220.23 & 354.01 & 464.02 & 560.98 \\
%0    & {N/A} & {N/A} & 98.73  & 147.65 & 214.39 & 347.87 & 457.41 & 553.97 \\
%5    & {N/A} & {N/A} & 88.85  & 141.28 & 208.56 & 341.66 & 450.70 & 546.83 \\
%10   & {N/A} & {N/A} & 78.39  & 135.03 & 202.80 & 335.46 & 443.96 & 539.61 \\
%15   & {N/A} & {N/A} & 68.14  & 129.05 & 197.20 & 329.34 & 437.24 & 532.38 \\
%20   & {N/A} & {N/A} & 59.71  & 123.57 & 191.88 & 323.36 & 430.63 & 525.22 \\
%25   & {N/A} & {N/A} & 54.06  & 118.69 & 186.91 & 317.60 & 424.16 & 518.19 \\
%30   & {N/A} & 27.24 & 51.06  & 114.50 & 182.34 & 312.10 & 417.94 & 511.36 \\
%35   & {N/A} & 26.78 & 50.12  & 111.06 & 178.22 & 306.91 & 412.01 & 504.81 \\
%40   & {N/A} & 27.70 & 50.34  & 108.50 & 174.68 & 302.14 & 406.39 & 498.55 \\
%45   & {N/A} & 29.08 & 51.39  & 106.67 & 171.68 & 297.78 & 401.14 & 492.63 \\
%50   & {N/A} & 30.74 & 52.90  & 105.51 & 169.21 & 293.85 & 396.35 & 487.18 \\
%55   & {N/A} & 32.52 & 54.76  & 104.96 & 167.26 & 290.38 & 392.01 & 482.16 \\
%60   & {N/A} & 34.39 & 56.87  & 104.98 & 165.89 & 287.44 & 388.13 & 477.60 \\
%65   & {N/A} & 36.30 & 59.12  & 105.46 & 165.00 & 284.97 & 384.75 & 473.51 \\
%70   & {N/A} & 38.25 & 61.48  & 106.30 & 164.55 & 282.97 & 381.85 & 469.98 \\
%75   & {N/A} & 40.20 & 63.90  & 107.45 & 164.50 & 281.41 & 379.42 & 466.92 \\
%80   & {N/A} & 42.15 & 66.38  & 108.91 & 164.86 & 280.36 & 377.53 & 464.42 \\
%85   & {N/A} & 44.09 & 68.87  & 110.55 & 165.54 & 279.72 & 376.11 & 462.41 \\
%90   & {N/A} & 46.02 & 71.37  & 112.40 & 166.54 & 279.47 & 375.11 & 460.87 \\
%95   & {N/A} & 47.92 & 73.86  & 114.37 & 167.78 & 279.58 & 374.54 & 459.79 \\
%100  & {N/A} & 49.80 & 76.32  & 116.45 & 169.24 & 280.07 & 374.41 & 459.18 \\
%105  & {N/A} & 51.65 & 78.76  & 118.62 & 170.87 & 280.86 & 374.65 & 458.98 \\
%110  & {N/A} & 53.46 & 81.15  & 120.84 & 172.69 & 281.91 & 375.21 & 459.21 \\
%115  & {N/A} & 55.23 & 83.49  & 123.09 & 174.61 & 283.20 & 376.08 & 459.77 \\
%120  & {N/A} & 56.96 & 85.77  & 125.35 & 176.62 & 284.72 & 377.25 & 460.65 \\
%125  & {N/A} & 58.64 & 88.00  & 127.60 & 178.70 & 286.41 & 378.65 & 461.81 \\
%130  & 48.33 & 60.28 & 90.14  & 129.82 & 180.82 & 288.22 & 380.23 & 463.22 \\
%135  & 48.34 & 61.86 & 92.20  & 132.00 & 182.95 & 290.14 & 381.98 & 464.84 \\
%140  & 48.55 & 63.38 & 94.18  & 134.12 & 185.06 & 292.14 & 383.87 & 466.63 \\
%145  & 48.89 & 64.83 & 96.07  & 136.17 & 187.15 & 294.18 & 385.84 & 468.54 \\
%150  & 49.33 & 66.23 & 97.85  & 138.13 & 189.18 & 296.23 & 387.87 & 470.53 \\
%155  & 49.84 & 67.57 & 99.54  & 139.99 & 191.14 & 298.25 & 389.91 & 472.58 \\
%160  & 50.39 & 68.83 & 101.12 & 141.74 & 193.00 & 300.22 & 391.94 & 474.63 \\
%165  & 50.97 & 70.03 & 102.60 & 143.39 & 194.76 & 302.13 & 393.92 & 476.66 \\
%170  & 51.58 & 71.15 & 103.97 & 144.91 & 196.41 & 303.94 & 395.83 & 478.65 \\
%175  & 52.20 & 72.21 & 105.23 & 146.30 & 197.92 & 305.65 & 397.64 & 480.54 \\
%180  & 52.81 & 73.20 & 106.38 & 147.58 & 199.32 & 307.22 & 399.34 & 482.33 \\
%185  & 53.43 & 74.12 & 107.42 & 148.73 & 200.58 & 308.66 & 400.93 & 484.00 \\
%190  & 54.04 & 74.97 & 108.37 & 149.75 & 201.70 & 309.97 & 402.35 & 485.55 \\
%195  & 54.64 & 75.75 & 109.21 & 150.66 & 202.69 & 311.14 & 403.64 & 486.93 \\
%200  & 55.23 & 76.47 & 109.95 & 151.45 & 203.55 & 312.15 & 404.78 & 488.18 \\
%205  & 55.80 & 77.13 & 110.60 & 152.13 & 204.30 & 313.04 & 405.79 & 489.27 \\
%210  & 56.35 & 77.72 & 111.16 & 152.70 & 204.93 & 313.80 & 406.65 & 490.23 \\
%215  & 56.88 & 78.25 & 111.65 & 153.19 & 205.45 & 314.44 & 407.39 & 491.06 \\
%220  & 57.39 & 78.72 & 112.04 & 153.58 & 205.87 & 314.97 & 408.01 & 491.77 \\
%225  & 57.87 & 79.14 & 112.38 & 153.91 & 206.21 & 315.39 & 408.53 & 492.35 \\
%230  & 58.34 & 79.51 & 112.64 & 154.14 & 206.47 & 315.73 & 408.94 & 492.83 \\
%235  & 58.78 & 79.83 & 112.85 & 154.31 & 206.66 & 316.00 & 409.27 & 493.22 \\
%240  & 59.20 & 80.10 & 113.00 & 154.43 & 206.79 & 316.18 & 409.52 & 493.54 \\
%245  & 59.60 & 80.33 & 113.10 & 154.49 & 206.86 & 316.30 & 409.70 & 493.78 \\
%250  & 59.97 & 80.50 & 113.15 & 154.51 & 206.88 & 316.38 & 409.84 & 493.96 \\
%255  & 60.32 & 80.65 & 113.16 & 154.49 & 206.86 & 316.41 & 409.92 & 494.09 \\
%260  & 60.65 & 80.76 & 113.13 & 154.43 & 206.80 & 316.41 & 409.96 & 494.18 \\
%265  & 60.97 & 80.83 & 113.07 & 154.34 & 206.73 & 316.37 & 409.97 & 494.24 \\
%270  & 61.25 & 80.86 & 112.99 & 154.23 & 206.63 & 316.31 & 409.96 & 494.27 \\
%\bottomrule
%\end{tabular}
%\end{table*}



\begin{table*}[tbp]
\centering
\footnotesize
\renewcommand{\arraystretch}{0.95}
\caption{Minimum transfer time (in days) from a circular orbit of 1 AU and final true anomaly (at rocket arrival).}
\label{tab:transfer_times_final_anomaly}
\sisetup{table-format=3.2} % Requires siunitx LaTeX package
\begin{tabular}{{r S S S S S S S }}
\toprule
\textbf{True Anomaly} & \multicolumn{7}{c}{\textbf{Target Radius (AU)}} \\
\cmidrule(lr){2-8}
\textbf{(\si{\degree})} & {0.390} & {0.720} & {1.520} & {5.200} & {9.580} & {19.220} & {30.050} \\
\midrule
0 & {N/A} & {N/A} & {N/A} & 220.50 & 337.82 & 539.55 & 724.65 \\
5 & {N/A} & {N/A} & 84.99 & 214.66 & 331.66 & 532.47 & 716.76 \\
10 & {N/A} & {N/A} & 73.01 & 208.91 & 325.52 & 525.31 & 708.75 \\
15 & {N/A} & {N/A} & 60.88 & 203.32 & 319.44 & 518.17 & 700.69 \\
20 & {N/A} & {N/A} & 51.35 & 197.98 & 313.53 & 511.09 & 692.64 \\
25 & {N/A} & {N/A} & 45.94 & 192.95 & 307.83 & 504.13 & 684.71 \\
30 & {N/A} & 32.66 & 43.80 & 188.35 & 302.40 & 497.40 & 676.93 \\
35 & {N/A} & 29.03 & 43.63 & 184.23 & 297.29 & 490.93 & 669.37 \\
40 & {N/A} & 28.92 & 44.61 & 180.59 & 292.59 & 484.76 & 662.08 \\
45 & {N/A} & 29.86 & 46.15 & 177.47 & 288.32 & 478.95 & 655.19 \\
50 & {N/A} & 31.19 & 48.07 & 174.94 & 284.48 & 473.58 & 648.68 \\
55 & {N/A} & 32.74 & 50.23 & 172.94 & 281.10 & 468.66 & 642.60 \\
60 & {N/A} & 34.42 & 52.55 & 171.44 & 278.25 & 464.20 & 637.00 \\
65 & {N/A} & 36.18 & 54.96 & 170.42 & 275.88 & 460.21 & 631.92 \\
70 & {N/A} & 38.00 & 57.46 & 169.90 & 273.97 & 456.76 & 627.37 \\
75 & {N/A} & 39.84 & 59.98 & 169.80 & 272.50 & 453.79 & 623.38 \\
80 & {N/A} & 41.69 & 62.51 & 170.05 & 271.54 & 451.39 & 619.93 \\
85 & {N/A} & 43.55 & 65.05 & 170.65 & 270.98 & 449.42 & 617.06 \\
90 & 46.97 & 45.39 & 67.58 & 171.57 & 270.81 & 447.98 & 614.72 \\
95 & 45.78 & 47.22 & 70.07 & 172.77 & 270.99 & 447.01 & 612.89 \\
100 & 45.47 & 49.04 & 72.55 & 174.17 & 271.54 & 446.46 & 611.58 \\
105 & 45.54 & 50.82 & 74.97 & 175.77 & 272.38 & 446.33 & 610.73 \\
110 & 45.86 & 52.58 & 77.35 & 177.53 & 273.50 & 446.57 & 610.37 \\
115 & 46.37 & 54.30 & 79.68 & 179.42 & 274.83 & 447.19 & 610.43 \\
120 & 47.00 & 55.99 & 81.93 & 181.41 & 276.38 & 448.11 & 610.89 \\
125 & 47.72 & 57.63 & 84.12 & 183.48 & 278.09 & 449.29 & 611.70 \\
130 & 48.49 & 59.22 & 86.24 & 185.58 & 279.93 & 450.73 & 612.82 \\
135 & 49.31 & 60.76 & 88.27 & 187.70 & 281.87 & 452.36 & 614.21 \\
140 & 50.16 & 62.24 & 90.22 & 189.81 & 283.87 & 454.17 & 615.83 \\
145 & 51.02 & 63.68 & 92.07 & 191.90 & 285.92 & 456.10 & 617.64 \\
150 & 51.89 & 65.05 & 93.83 & 193.93 & 287.97 & 458.09 & 619.57 \\
155 & 52.75 & 66.36 & 95.50 & 195.89 & 289.99 & 460.14 & 621.61 \\
160 & 53.61 & 67.61 & 97.05 & 197.76 & 291.96 & 462.19 & 623.69 \\
165 & 54.46 & 68.79 & 98.50 & 199.53 & 293.86 & 464.22 & 625.79 \\
170 & 55.28 & 69.89 & 99.84 & 201.18 & 295.66 & 466.18 & 627.86 \\
175 & 56.08 & 70.94 & 101.08 & 202.72 & 297.35 & 468.07 & 629.88 \\
180 & 56.86 & 71.91 & 102.22 & 204.12 & 298.92 & 469.85 & 631.81 \\
185 & 57.61 & 72.83 & 103.25 & 205.38 & 300.34 & 471.50 & 633.63 \\
190 & 58.34 & 73.67 & 104.18 & 206.51 & 301.63 & 473.03 & 635.31 \\
195 & 59.03 & 74.46 & 105.02 & 207.53 & 302.79 & 474.41 & 636.87 \\
200 & 59.68 & 75.17 & 105.75 & 208.39 & 303.79 & 475.63 & 638.27 \\
205 & 60.31 & 75.83 & 106.40 & 209.13 & 304.67 & 476.72 & 639.53 \\
210 & 60.90 & 76.42 & 106.97 & 209.77 & 305.42 & 477.67 & 640.64 \\
215 & 61.45 & 76.96 & 107.46 & 210.31 & 306.05 & 478.48 & 641.61 \\
220 & 61.97 & 77.44 & 107.86 & 210.74 & 306.57 & 479.17 & 642.45 \\
225 & 62.47 & 77.86 & 108.20 & 211.07 & 306.99 & 479.74 & 643.16 \\
230 & 62.93 & 78.23 & 108.48 & 211.33 & 307.32 & 480.22 & 643.76 \\
235 & 63.35 & 78.56 & 108.69 & 211.52 & 307.57 & 480.60 & 644.27 \\
240 & 63.74 & 78.84 & 108.85 & 211.65 & 307.75 & 480.91 & 644.68 \\
245 & 64.10 & 79.07 & 108.95 & 211.72 & 307.87 & 481.14 & 645.02 \\
250 & 64.44 & 79.26 & 109.01 & 211.75 & 307.94 & 481.31 & 645.29 \\
255 & 64.74 & 79.42 & 109.03 & 211.73 & 307.97 & 481.43 & 645.51 \\
260 & 65.02 & 79.54 & 109.01 & 211.68 & 307.96 & 481.51 & 645.68 \\
265 & 65.26 & 79.62 & 108.96 & 211.60 & 307.92 & 481.57 & 645.82 \\
270 & 65.48 & 79.66 & 108.88 & 211.50 & 307.87 & 481.59 & 645.94 \\
275 & 65.69 & 79.69 & 108.78 & 211.40 & 307.79 & 481.61 & 646.04 \\
280 & 65.86 & 79.69 & 108.66 & 211.27 & 307.72 & 481.62 & 646.12 \\
285 & 66.02 & 79.65 & 108.52 & 211.14 & 307.63 & 481.62 & 646.22 \\
290 & 66.16 & 79.61 & 108.37 & 211.00 & 307.55 & 481.64 & 646.31 \\
295 & 66.27 & 79.54 & 108.21 & 210.88 & 307.48 & 481.67 & 646.43 \\
300 & 66.38 & 79.47 & 108.03 & 210.75 & 307.42 & 481.71 & 646.57 \\
305 & 66.47 & 79.37 & 107.87 & 210.65 & 307.38 & 481.79 & 646.74 \\
310 & 66.56 & 79.27 & 107.70 & 210.56 & 307.37 & 481.89 & 646.96 \\
315 & 66.63 & 79.18 & 107.54 & 210.49 & 307.38 & 482.04 & 647.22 \\
320 & 66.70 & 79.07 & 107.39 & 210.45 & 307.43 & 482.23 & 647.54 \\
325 & 66.77 & 78.98 & 107.26 & 210.43 & 307.52 & 482.48 & 647.92 \\
330 & 66.84 & 78.88 & 107.14 & 210.45 & 307.64 & 482.79 & 648.38 \\
335 & 66.90 & 78.81 & 107.03 & 210.51 & 307.82 & 483.16 & 648.91 \\
340 & 66.99 & 78.74 & 106.97 & 210.61 & 308.06 & 483.61 & 649.54 \\
345 & 67.08 & 78.70 & 106.94 & 210.75 & 308.36 & 484.14 & 650.27 \\
350 & 67.19 & 78.70 & 106.94 & 210.96 & 308.71 & 484.76 & 651.11 \\
355 & 67.32 & 78.73 & 107.00 & 211.22 & 309.15 & 485.47 & 652.06 \\
360 & 67.49 & 78.81 & 107.10 & 211.54 & 309.66 & 486.29 & 653.13 \\
\bottomrule
\end{tabular}
\end{table*}


\begin{table*}[tbp]
\centering
\footnotesize
\renewcommand{\arraystretch}{0.95}
\caption{Minimum transfer time (in days) from a circular orbit of 1 AU and initial true anomaly (at rocket departure).}
\label{tab:transfer_times_initial_anomaly}
\sisetup{table-format=3.2} % Requires siunitx LaTeX package
\begin{tabular}{{r S S S S S S S }}
\toprule
\textbf{True Anomaly} & \multicolumn{7}{c}{\textbf{Target Radius (AU)}} \\
\cmidrule(lr){2-8}
\textbf{(\si{\degree})} & {0.390} & {0.720} & {1.520} & {5.200} & {9.580} & {19.220} & {30.050} \\
\midrule
-100 & 46.96 & {N/A} & {N/A} & 281.19 & 401.02 & 604.63 & {N/A} \\
-95 & 46.26 & {N/A} & {N/A} & 279.72 & 400.45 & 605.12 & {N/A} \\
-90 & 45.76 & {N/A} & {N/A} & 277.93 & 399.55 & 605.20 & {N/A} \\
-85 & 45.48 & {N/A} & {N/A} & 275.79 & 398.22 & 604.85 & {N/A} \\
-80 & 45.50 & {N/A} & {N/A} & 273.34 & 396.54 & 604.07 & {N/A} \\
-75 & 45.92 & {N/A} & {N/A} & 270.57 & 394.49 & 602.85 & {N/A} \\
-70 & 47.06 & {N/A} & {N/A} & 267.50 & 392.06 & 601.12 & {N/A} \\
-65 & 49.08 & {N/A} & {N/A} & 264.11 & 389.20 & 598.91 & {N/A} \\
-60 & 51.87 & {N/A} & {N/A} & 260.40 & 385.96 & 596.20 & {N/A} \\
-55 & 54.65 & {N/A} & {N/A} & 256.39 & 382.34 & 592.99 & 783.40 \\
-50 & 56.90 & {N/A} & {N/A} & 252.12 & 378.37 & 589.31 & 779.93 \\
-45 & 58.68 & {N/A} & {N/A} & 247.60 & 374.05 & 585.12 & 775.84 \\
-40 & 60.08 & {N/A} & 85.30 & 242.86 & 369.37 & 580.49 & 771.21 \\
-35 & 61.19 & {N/A} & 80.06 & 237.92 & 364.40 & 575.39 & 766.05 \\
-30 & 62.11 & 36.38 & 74.72 & 232.82 & 359.17 & 569.87 & 760.36 \\
-25 & 62.88 & 33.69 & 69.37 & 227.58 & 353.71 & 564.00 & 754.17 \\
-20 & 63.51 & 31.65 & 64.02 & 222.28 & 348.06 & 557.81 & 747.56 \\
-15 & 64.05 & 29.77 & 58.87 & 216.96 & 342.25 & 551.34 & 740.58 \\
-10 & 64.50 & 28.86 & 53.88 & 211.67 & 336.37 & 544.62 & 733.26 \\
-5 & 64.89 & 29.30 & 49.74 & 206.46 & 330.45 & 537.74 & 725.67 \\
0 & 65.21 & 31.38 & 46.38 & 201.41 & 324.55 & 530.74 & 717.88 \\
5 & 65.48 & 34.85 & 44.28 & 196.58 & 318.72 & 523.70 & 709.96 \\
10 & 65.72 & 39.16 & 43.59 & 192.04 & 313.06 & 516.69 & 701.98 \\
15 & 65.92 & 43.78 & 44.17 & 187.83 & 307.58 & 509.73 & 694.01 \\
20 & 66.09 & 48.26 & 45.84 & 184.00 & 302.34 & 502.93 & 686.11 \\
25 & 66.23 & 52.44 & 48.21 & 180.58 & 297.40 & 496.34 & 678.37 \\
30 & 66.35 & 56.21 & 51.07 & 177.61 & 292.82 & 490.00 & 670.84 \\
35 & 66.45 & 59.57 & 54.25 & 175.13 & 288.65 & 483.96 & 663.57 \\
40 & 66.55 & 62.47 & 57.60 & 173.15 & 284.87 & 478.29 & 656.62 \\
45 & 66.63 & 65.02 & 61.04 & 171.63 & 281.52 & 473.04 & 650.08 \\
50 & 66.70 & 67.23 & 64.51 & 170.57 & 278.62 & 468.22 & 643.97 \\
55 & 66.77 & 69.13 & 67.95 & 169.94 & 276.22 & 463.83 & 638.30 \\
60 & 66.84 & 70.78 & 71.31 & 169.78 & 274.27 & 459.92 & 633.10 \\
65 & 66.91 & 72.19 & 74.60 & 169.99 & 272.76 & 456.56 & 628.48 \\
70 & 67.01 & 73.43 & 77.78 & 170.54 & 271.68 & 453.63 & 624.33 \\
75 & 67.11 & 74.48 & 80.81 & 171.40 & 271.05 & 451.26 & 620.80 \\
80 & 67.24 & 75.39 & 83.70 & 172.59 & 270.82 & 449.33 & 617.74 \\
85 & 67.41 & 76.18 & 86.43 & 174.01 & 270.94 & 447.93 & 615.27 \\
90 & 67.62 & 76.85 & 88.98 & 175.63 & 271.40 & 446.97 & 613.30 \\
95 & 67.90 & 77.44 & 91.37 & 177.43 & 272.19 & 446.44 & 611.88 \\
100 & 68.28 & 77.91 & 93.55 & 179.39 & 273.28 & 446.33 & 610.94 \\
105 & 68.79 & 78.33 & 95.58 & 181.45 & 274.59 & 446.59 & 610.43 \\
110 & 69.52 & 78.67 & 97.41 & 183.58 & 276.10 & 447.23 & 610.40 \\
115 & {N/A} & 78.96 & 99.09 & 185.77 & 277.82 & 448.16 & 610.75 \\
120 & {N/A} & 79.18 & 100.57 & 187.96 & 279.67 & 449.36 & 611.45 \\
125 & {N/A} & 79.36 & 101.92 & 190.15 & 281.62 & 450.81 & 612.49 \\
130 & {N/A} & 79.50 & 103.10 & 192.30 & 283.64 & 452.46 & 613.81 \\
135 & {N/A} & 79.59 & 104.14 & 194.39 & 285.71 & 454.29 & 615.39 \\
140 & {N/A} & 79.65 & 105.06 & 196.39 & 287.79 & 456.23 & 617.15 \\
145 & {N/A} & 79.68 & 105.84 & 198.29 & 289.84 & 458.24 & 619.06 \\
150 & {N/A} & 79.69 & 106.53 & 200.09 & 291.84 & 460.29 & 621.07 \\
155 & {N/A} & 79.66 & 107.11 & 201.74 & 293.77 & 462.35 & 623.15 \\
160 & {N/A} & 79.63 & 107.59 & 203.26 & 295.60 & 464.39 & 625.26 \\
165 & {N/A} & 79.57 & 107.99 & 204.64 & 297.31 & 466.36 & 627.35 \\
170 & {N/A} & 79.49 & 108.32 & 205.88 & 298.90 & 468.25 & 629.38 \\
175 & {N/A} & 79.41 & 108.58 & 206.97 & 300.33 & 470.02 & 631.35 \\
180 & {N/A} & 79.32 & 108.76 & 207.93 & 301.64 & 471.67 & 633.19 \\
185 & {N/A} & 79.22 & 108.90 & 208.75 & 302.80 & 473.19 & 634.91 \\
190 & {N/A} & 79.12 & 108.99 & 209.46 & 303.81 & 474.55 & 636.50 \\
195 & {N/A} & 79.02 & 109.02 & 210.05 & 304.69 & 475.76 & 637.95 \\
200 & {N/A} & 78.93 & 109.02 & 210.53 & 305.44 & 476.84 & 639.25 \\
205 & {N/A} & 78.85 & 109.00 & 210.92 & 306.07 & 477.78 & 640.39 \\
210 & {N/A} & 78.78 & 108.93 & 211.22 & 306.59 & 478.57 & 641.39 \\
215 & {N/A} & 78.72 & 108.84 & 211.44 & 307.00 & 479.25 & 642.26 \\
220 & {N/A} & 78.70 & 108.73 & 211.59 & 307.33 & 479.80 & 643.01 \\
225 & {N/A} & 78.71 & 108.60 & 211.69 & 307.58 & 480.28 & 643.63 \\
230 & {N/A} & 78.76 & 108.46 & 211.74 & 307.76 & 480.64 & 644.16 \\
235 & {N/A} & 78.87 & 108.31 & 211.74 & 307.88 & 480.94 & 644.60 \\
240 & {N/A} & 79.06 & 108.14 & 211.71 & 307.95 & 481.16 & 644.95 \\
245 & {N/A} & 79.33 & 107.98 & 211.64 & 307.97 & 481.33 & 645.23 \\
250 & {N/A} & 79.74 & 107.81 & 211.55 & 307.96 & 481.44 & 645.46 \\
255 & {N/A} & 80.32 & 107.65 & 211.45 & 307.92 & 481.52 & 645.64 \\
260 & {N/A} & 81.16 & 107.49 & 211.33 & 307.86 & 481.57 & 645.79 \\
\bottomrule
\end{tabular}
\end{table*}



\begin{table*}[!htbp]
\centering
\caption{Minimum transfer time (in days) to a circular orbit of 1 AU and final true anomaly (at rocket arrival).}
\label{tab:transfer_times_final_anomaly_return}
\sisetup{table-format=3.2} % Requires siunitx LaTeX package
\begin{tabular}{{r S S S S S S S }}
\toprule
\textbf{True Anomaly} & \multicolumn{7}{c}{\textbf{Origin Radius (AU)}} \\
\cmidrule(lr){2-8}
\textbf{(\si{\degree})} & {0.390} & {0.720} & {1.520} & {5.200} & {9.580} & {19.220} & {30.050} \\
\midrule
0 & {N/A} & {N/A} & 95.77 & 220.50 & 337.82 & 539.56 & 724.64 \\
5 & {N/A} & {N/A} & 84.93 & 214.67 & 331.66 & 532.46 & 716.76 \\
10 & {N/A} & {N/A} & 73.02 & 208.91 & 325.51 & 525.31 & 708.76 \\
15 & {N/A} & {N/A} & 60.90 & 203.30 & 319.44 & 518.16 & 700.69 \\
20 & {N/A} & {N/A} & 51.21 & 197.97 & 313.53 & 511.08 & 692.65 \\
25 & {N/A} & {N/A} & 45.92 & 192.97 & 307.82 & 504.14 & 684.70 \\
30 & {N/A} & 32.29 & 43.85 & 188.36 & 302.39 & 497.40 & 676.91 \\
35 & {N/A} & 28.90 & 43.65 & 184.19 & 297.29 & 490.93 & 669.34 \\
40 & {N/A} & 28.92 & 44.59 & 180.58 & 292.60 & 484.76 & 662.08 \\
45 & {N/A} & 29.84 & 46.15 & 177.50 & 288.32 & 478.95 & 655.19 \\
50 & {N/A} & 31.18 & 48.07 & 174.95 & 284.47 & 473.58 & 648.68 \\
55 & {N/A} & 32.73 & 50.23 & 172.91 & 281.10 & 468.66 & 642.61 \\
60 & {N/A} & 34.42 & 52.55 & 171.44 & 278.26 & 464.19 & 637.03 \\
65 & {N/A} & 36.19 & 54.97 & 170.45 & 275.88 & 460.24 & 631.95 \\
70 & {N/A} & 37.99 & 57.46 & 169.91 & 273.97 & 456.77 & 627.37 \\
75 & {N/A} & 39.83 & 59.98 & 169.77 & 272.50 & 453.80 & 623.40 \\
80 & {N/A} & 41.69 & 62.51 & 170.05 & 271.54 & 451.38 & 619.96 \\
85 & 49.67 & 43.54 & 65.05 & 170.67 & 270.98 & 449.42 & 617.07 \\
90 & 46.93 & 45.39 & 67.57 & 171.58 & 270.81 & 447.98 & 614.71 \\
95 & 45.83 & 47.22 & 70.08 & 172.74 & 270.99 & 447.01 & 612.86 \\
100 & 45.46 & 49.03 & 72.54 & 174.16 & 271.54 & 446.45 & 611.58 \\
105 & 45.51 & 50.82 & 74.97 & 175.77 & 272.38 & 446.33 & 610.73 \\
110 & 45.85 & 52.58 & 77.36 & 177.54 & 273.50 & 446.58 & 610.38 \\
115 & 46.37 & 54.31 & 79.67 & 179.42 & 274.83 & 447.17 & 610.46 \\
120 & 47.00 & 55.99 & 81.93 & 181.41 & 276.38 & 448.10 & 610.91 \\
125 & 47.72 & 57.63 & 84.12 & 183.47 & 278.09 & 449.30 & 611.70 \\
130 & 48.49 & 59.22 & 86.23 & 185.58 & 279.93 & 450.74 & 612.82 \\
135 & 49.31 & 60.76 & 88.27 & 187.70 & 281.87 & 452.36 & 614.21 \\
140 & 50.15 & 62.24 & 90.22 & 189.81 & 283.87 & 454.17 & 615.83 \\
145 & 51.02 & 63.68 & 92.07 & 191.90 & 285.92 & 456.10 & 617.64 \\
150 & 51.89 & 65.04 & 93.84 & 193.93 & 287.97 & 458.09 & 619.58 \\
155 & 52.75 & 66.35 & 95.49 & 195.89 & 289.99 & 460.14 & 621.61 \\
160 & 53.61 & 67.60 & 97.05 & 197.76 & 291.96 & 462.19 & 623.69 \\
165 & 54.45 & 68.79 & 98.49 & 199.53 & 293.86 & 464.22 & 625.79 \\
170 & 55.28 & 69.89 & 99.84 & 201.18 & 295.66 & 466.18 & 627.86 \\
175 & 56.08 & 70.94 & 101.08 & 202.72 & 297.35 & 468.07 & 629.88 \\
180 & 56.86 & 71.91 & 102.23 & 204.13 & 298.92 & 469.85 & 631.81 \\
185 & 57.61 & 72.83 & 103.26 & 205.39 & 300.35 & 471.50 & 633.62 \\
190 & 58.33 & 73.68 & 104.18 & 206.52 & 301.64 & 473.03 & 635.31 \\
195 & 59.02 & 74.46 & 105.02 & 207.52 & 302.79 & 474.40 & 636.87 \\
200 & 59.68 & 75.17 & 105.76 & 208.39 & 303.80 & 475.63 & 638.27 \\
205 & 60.31 & 75.83 & 106.41 & 209.14 & 304.67 & 476.72 & 639.52 \\
210 & 60.90 & 76.42 & 106.97 & 209.77 & 305.42 & 477.67 & 640.64 \\
215 & 61.46 & 76.96 & 107.45 & 210.30 & 306.06 & 478.48 & 641.61 \\
220 & 61.98 & 77.44 & 107.87 & 210.73 & 306.58 & 479.17 & 642.44 \\
225 & 62.47 & 77.86 & 108.21 & 211.07 & 306.99 & 479.75 & 643.17 \\
230 & 62.93 & 78.23 & 108.48 & 211.34 & 307.32 & 480.22 & 643.76 \\
235 & 63.35 & 78.57 & 108.69 & 211.52 & 307.58 & 480.60 & 644.27 \\
240 & 63.75 & 78.84 & 108.84 & 211.65 & 307.75 & 480.90 & 644.68 \\
245 & 64.10 & 79.07 & 108.95 & 211.72 & 307.87 & 481.14 & 645.02 \\
250 & 64.44 & 79.27 & 109.01 & 211.75 & 307.94 & 481.31 & 645.30 \\
255 & 64.74 & 79.41 & 109.03 & 211.73 & 307.97 & 481.43 & 645.51 \\
260 & 65.01 & 79.53 & 109.01 & 211.68 & 307.96 & 481.51 & 645.68 \\
265 & 65.26 & 79.62 & 108.97 & 211.61 & 307.92 & 481.57 & 645.82 \\
270 & 65.49 & 79.66 & 108.88 & 211.51 & 307.87 & 481.60 & 645.93 \\
275 & 65.69 & 79.68 & 108.78 & 211.39 & {N/A} & 481.61 & 646.03 \\
280 & 65.86 & 79.68 & 108.66 & 211.26 & {N/A} & 481.62 & 646.12 \\
285 & 66.02 & 79.65 & 108.52 & 211.14 & {N/A} & 481.62 & 646.21 \\
290 & 66.16 & 79.61 & 108.36 & 211.00 & {N/A} & 481.63 & 646.31 \\
295 & 66.27 & 79.54 & 108.21 & 210.88 & {N/A} & 481.66 & 646.43 \\
300 & 66.38 & 79.47 & 108.04 & 210.75 & {N/A} & 481.71 & 646.57 \\
305 & 66.47 & 79.37 & 107.87 & 210.65 & {N/A} & 481.78 & 646.75 \\
310 & 66.56 & 79.27 & 107.70 & 210.56 & {N/A} & 481.89 & 646.96 \\
315 & 66.63 & 79.17 & 107.54 & 210.49 & {N/A} & 482.04 & 647.22 \\
320 & 66.70 & 79.07 & 107.39 & 210.44 & {N/A} & 482.24 & 647.54 \\
325 & 66.76 & 78.97 & 107.25 & 210.43 & {N/A} & 482.48 & 647.93 \\
330 & 66.83 & 78.89 & 107.13 & 210.45 & {N/A} & 482.79 & 648.38 \\
335 & 66.90 & 78.80 & 107.03 & 210.51 & {N/A} & 483.16 & 648.91 \\
340 & 66.99 & 78.74 & 106.97 & 210.60 & {N/A} & 483.61 & 649.55 \\
345 & 67.08 & 78.70 & 106.94 & 210.75 & {N/A} & 484.14 & 650.27 \\
350 & 67.19 & 78.70 & 106.94 & 210.96 & {N/A} & 484.76 & 651.11 \\
355 & 67.32 & 78.73 & 107.00 & 211.22 & {N/A} & 485.47 & 652.06 \\
360 & 67.48 & 78.82 & 107.10 & 211.54 & {N/A} & 486.28 & 653.13 \\
\bottomrule
\end{tabular}
\end{table*}


\begin{table*}[!htbp]
\centering
\caption{Minimum transfer time (in days) to a circular orbit of 1 AU and initial true anomaly (at rocket departure).}
\label{tab:transfer_times_initial_anomaly_return}
\sisetup{table-format=3.2} % Requires siunitx LaTeX package
\begin{tabular}{{r S S S S S S S }}
\toprule
\textbf{True Anomaly} & \multicolumn{7}{c}{\textbf{Origin Radius (AU)}} \\
\cmidrule(lr){2-8}
\textbf{(\si{\degree})} & {0.390} & {0.720} & {1.520} & {5.200} & {9.580} & {19.220} & {30.050} \\
\midrule
-100 & 46.94 & {N/A} & 137.22 & 281.20 & 401.02 & 604.64 & 789.45 \\
-95 & 46.28 & {N/A} & 133.71 & 279.71 & 400.45 & 605.10 & 790.72 \\
-90 & 45.79 & {N/A} & 130.04 & 277.91 & 399.53 & 605.19 & 791.59 \\
-85 & 45.47 & {N/A} & 126.23 & 275.81 & 398.24 & 604.86 & 792.02 \\
-80 & 45.48 & {N/A} & 122.29 & 273.35 & 396.55 & 604.09 & 791.99 \\
-75 & 45.94 & {N/A} & 118.16 & 270.58 & 394.50 & 602.83 & 791.40 \\
-70 & 47.06 & {N/A} & 113.88 & 267.49 & 392.03 & 601.11 & 790.28 \\
-65 & 49.09 & {N/A} & 109.49 & 264.09 & 389.19 & 598.93 & 788.57 \\
-60 & 51.87 & {N/A} & 104.93 & 260.40 & 385.96 & 596.20 & 786.30 \\
-55 & 54.64 & {N/A} & 100.21 & 256.39 & 382.37 & 593.01 & 783.38 \\
-50 & 56.91 & {N/A} & 95.40 & 252.12 & 378.37 & 589.32 & 779.89 \\
-45 & 58.67 & {N/A} & 90.40 & 247.61 & 374.03 & 585.12 & 775.84 \\
-40 & 60.08 & 42.22 & 85.27 & 242.85 & 369.36 & 580.46 & 771.21 \\
-35 & 61.20 & 39.24 & 80.06 & 237.91 & 364.41 & 575.38 & 766.07 \\
-30 & 62.11 & 36.39 & 74.72 & 232.81 & 359.17 & 569.87 & 760.34 \\
-25 & 62.88 & 33.74 & 69.36 & 227.59 & 353.71 & 564.01 & 754.16 \\
-20 & 63.51 & 31.37 & 64.05 & 222.28 & 348.05 & 557.81 & 747.56 \\
-15 & 64.05 & 29.82 & 58.75 & 216.95 & 342.25 & 551.33 & 740.58 \\
-10 & 64.50 & 28.85 & 53.84 & 211.66 & 336.37 & 544.62 & 733.26 \\
-5 & 64.88 & 29.29 & 49.74 & 206.47 & 330.45 & 537.73 & 725.67 \\
0 & 65.21 & 31.36 & 46.31 & 201.42 & 324.55 & 530.74 & 717.89 \\
5 & 65.49 & 34.83 & 44.22 & 196.60 & 318.73 & 523.70 & 709.96 \\
10 & 65.72 & 39.16 & 43.58 & 192.04 & 313.06 & 516.67 & 701.98 \\
15 & 65.92 & 43.77 & 44.15 & 187.81 & 307.57 & 509.73 & 694.00 \\
20 & 66.09 & 48.26 & 45.83 & 183.96 & 302.33 & 502.94 & 686.10 \\
25 & 66.23 & 52.45 & 48.22 & 180.57 & 297.40 & 496.34 & 678.35 \\
30 & 66.35 & 56.21 & 51.07 & 177.63 & 292.84 & 490.00 & 670.82 \\
35 & 66.46 & 59.57 & 54.25 & 175.16 & 288.65 & 483.95 & 663.55 \\
40 & 66.54 & 62.47 & 57.61 & 173.14 & 284.87 & 478.29 & 656.62 \\
45 & 66.63 & 65.02 & 61.04 & 171.60 & 281.51 & 473.04 & 650.08 \\
50 & 66.70 & 67.22 & 64.50 & 170.57 & 278.64 & 468.22 & 643.97 \\
55 & 66.77 & 69.13 & 67.94 & 169.97 & 276.22 & 463.84 & 638.28 \\
60 & 66.84 & 70.77 & 71.32 & 169.78 & 274.27 & 459.96 & 633.13 \\
65 & 66.92 & 72.19 & 74.61 & 169.97 & 272.76 & 456.55 & 628.48 \\
70 & 67.01 & 73.42 & 77.78 & 170.54 & 271.68 & 453.63 & 624.35 \\
75 & 67.12 & 74.48 & 80.81 & 171.43 & 271.05 & 451.26 & 620.77 \\
80 & 67.24 & 75.39 & 83.70 & 172.59 & 270.82 & 449.33 & 617.75 \\
85 & 67.41 & 76.18 & 86.42 & 173.99 & 270.94 & 447.93 & 615.28 \\
90 & 67.62 & 76.85 & 88.98 & 175.64 & 271.40 & 446.97 & 613.33 \\
95 & {N/A} & 77.43 & 91.36 & 177.44 & 272.19 & 446.44 & 611.88 \\
100 & {N/A} & 77.92 & 93.56 & 179.39 & 273.28 & 446.33 & 610.94 \\
105 & {N/A} & 78.33 & 95.57 & 181.45 & 274.59 & 446.60 & 610.43 \\
110 & {N/A} & 78.67 & 97.42 & 183.58 & 276.10 & 447.20 & 610.38 \\
115 & {N/A} & 78.95 & 99.07 & 185.77 & 277.82 & 448.15 & 610.75 \\
120 & {N/A} & 79.18 & 100.57 & 187.97 & 279.67 & 449.37 & 611.47 \\
125 & {N/A} & 79.36 & 101.92 & 190.15 & 281.62 & 450.82 & 612.50 \\
130 & {N/A} & 79.49 & 103.11 & 192.30 & 283.64 & 452.46 & 613.81 \\
135 & {N/A} & 79.59 & 104.14 & 194.39 & 285.71 & 454.29 & 615.39 \\
140 & {N/A} & 79.65 & 105.06 & 196.39 & 287.79 & 456.23 & 617.15 \\
145 & {N/A} & 79.68 & 105.85 & 198.29 & 289.84 & 458.24 & 619.06 \\
150 & {N/A} & 79.69 & 106.54 & 200.09 & 291.84 & 460.29 & 621.08 \\
155 & {N/A} & 79.66 & 107.11 & 201.74 & 293.77 & 462.35 & 623.15 \\
160 & {N/A} & 79.63 & 107.60 & 203.26 & 295.60 & 464.39 & 625.26 \\
165 & {N/A} & 79.57 & 108.00 & 204.64 & 297.31 & 466.36 & 627.35 \\
170 & {N/A} & 79.49 & 108.32 & 205.88 & 298.90 & 468.25 & 629.38 \\
175 & {N/A} & 79.41 & 108.58 & 206.99 & 300.34 & 470.02 & 631.35 \\
180 & {N/A} & 79.32 & 108.76 & 207.93 & 301.65 & 471.67 & 633.20 \\
185 & {N/A} & 79.22 & 108.90 & 208.75 & 302.80 & 473.19 & 634.91 \\
190 & {N/A} & 79.12 & 108.98 & 209.46 & 303.82 & 474.54 & 636.50 \\
195 & {N/A} & 79.02 & 109.02 & 210.05 & 304.69 & 475.76 & 637.95 \\
200 & {N/A} & 78.93 & 109.03 & 210.52 & 305.45 & 476.84 & 639.25 \\
205 & {N/A} & 78.85 & 108.99 & 210.92 & 306.08 & 477.77 & 640.40 \\
210 & {N/A} & 78.77 & 108.93 & 211.22 & 306.59 & 478.57 & 641.40 \\
215 & {N/A} & 78.72 & 108.84 & 211.44 & 307.01 & 479.25 & 642.26 \\
220 & {N/A} & 78.70 & 108.73 & 211.59 & 307.34 & 479.81 & 643.01 \\
225 & {N/A} & 78.71 & 108.60 & 211.70 & 307.59 & 480.28 & 643.64 \\
230 & {N/A} & 78.76 & 108.46 & 211.74 & 307.76 & 480.64 & 644.17 \\
235 & {N/A} & 78.87 & 108.30 & 211.74 & 307.88 & 480.93 & 644.60 \\
240 & {N/A} & 79.06 & 108.15 & 211.71 & 307.94 & 481.17 & 644.95 \\
245 & {N/A} & {N/A} & 107.97 & 211.64 & 307.97 & 481.33 & 645.24 \\
250 & {N/A} & {N/A} & 107.81 & 211.56 & 307.96 & 481.44 & 645.47 \\
255 & {N/A} & {N/A} & 107.65 & 211.45 & 307.92 & 481.52 & 645.64 \\
260 & {N/A} & {N/A} & 107.49 & 211.33 & 307.87 & 481.57 & 645.79 \\
\bottomrule
\end{tabular}
\end{table*}



\end{document}